\documentclass[11pt]{article}
\usepackage[margin=1in]{geometry}
\usepackage{amsmath,amssymb,amsthm,mathtools}
\usepackage[T1]{fontenc}
\usepackage{microtype}
\usepackage{booktabs}
\usepackage{longtable}
\usepackage{enumitem}
\usepackage{xcolor}
\usepackage{hyperref}
\hypersetup{colorlinks=true,linkcolor=blue!50!black,citecolor=blue!50!black,urlcolor=blue!50!black}

\theoremstyle{plain}
\newtheorem{theorem}{Theorem}[section]
\newtheorem{proposition}[theorem]{Proposition}
\newtheorem{lemma}[theorem]{Lemma}
\newtheorem{corollary}[theorem]{Corollary}
\theoremstyle{definition}
\newtheorem{definition}[theorem]{Definition}
\newtheorem{conjecture}[theorem]{Conjecture}
\newtheorem{openproblem}{Open Problem}
\renewcommand{\theopenproblem}{\Alph{openproblem}}
\theoremstyle{remark}
\newtheorem{remark}[theorem]{Remark}
\newtheorem*{remark*}{Remark}
\newtheorem{observation}[theorem]{Observation}

\newcommand{\Occ}{\mathcal{O}}
\newcommand{\Edep}{\mathsf{E}}          % endpoint depth, distinct from carry excess E_n
\newcommand{\Zt}{\mathbb{Z}_2}
\newcommand{\Zpos}{\mathbb{Z}_{>0}}
\newcommand{\al}{\alpha}
\newcommand{\rmin}{r_{\min}}
\newcommand{\status}[1]{\textsf{\footnotesize[#1]}}
\newcommand{\doi}[1]{\textsf{\footnotesize doi:}\href{https://doi.org/#1}{\texttt{\footnotesize #1}}}

\title{A Global Occupation Conjecture for the Accelerated $3x+1$ Map:\\
Residue--Time Duality, Divergent-Orbit Sparsity,\\ and the Pointwise Realizer Problem}
\author{Elias De Jes\'us\\ Independent Researcher\\ ORCID: 0009-0007-0190-9143\\ \texttt{dejesuselias10@gmail.com}}
\date{Revision 7, 2026}

\begin{document}
\maketitle

\begin{abstract}
We study the Effective Occupation Conjecture (EOC) for the accelerated $3x+1$ map
$T(m)=(3m+1)/2^{a(m)}$, $a(m)=\nu_2(3m+1)$. Writing $R_n=S_n-n\log_2 3$ for the drift of the
accumulated $2$-adic valuation and $\Occ_c(m_0)=\#\{n\ge 0: R_n\le c\}$, EOC asserts
$\Occ_c(m_0)=O(\log m_0)$ for every seed and every corridor width~$c$.

This revision reorganizes the programme around four axes --- occupation, residue, time, and
divergence --- and the exact bridges between them. Record chronology identifies the residue and
time axes precisely: the least-realizer frontier $\rmin(N,c)$ is the inverse function of the
single-window confinement lifetime $L_c(m)$, so every pointwise lifetime bound transfers to a
realizer floor, on a graded scale that we tabulate. On the divergence side we retain the strongest
unconditional input: via Garcia--Tal's collision-free window theorem and Curry's explicit
quantitative form, every divergent orbit has summable reciprocals and admits no eventual
logarithmic drift floor $R_n\ge -B\log_2 n+O(1)$ for any $B<1/\beta_*\approx 1.0358567$. Curry's
input plays a second, qualitative role: after the last global maximum of $R$, a hypothetical
divergent orbit has an infinite zero-corridor tail, so \emph{universal drift exit} --- every
positive seed eventually attains $R_n>0$ --- is equivalent to the exclusion of divergent orbits,
while not excluding hypothetical nontrivial cycles, so that a proof of it alone would not yield the
Collatz conjecture --- no implication from it to cycle-freeness is established here
(Rem.~\ref{rem:strictlyweaker}).

We then record what a sustained programme of structural audits has closed. Finite information
cannot decide infinite realization: every exact-realizer cylinder contains arbitrarily large
positive integers, so $\forall N\,\exists m_N$ is free and all content lies in $\exists m\,\forall N$.
Nor do the \emph{known divergent-orbit} valuation-word shadows decide it: we exhibit an explicitly
indexed
uncountable family of zero-confined words with linear negative drift, summable $\sum 2^{R_n}$ and
$O(G)$ occupation shadow, of which all but countably many have no positive-integer realizer. Word-level
criteria in general are not excluded, and one is exact (Rem.~\ref{rem:zcrecounter}).
Several apparently independent coordinates --- normalized carry, lift digits, terminal zero gap,
transport budgets --- are invertible reparameterizations of the exact realizer, and the Archimedean
and $2$-adic readings of the natural divisibility object are complementary factors of one identity
rather than independent constraints. Launch renormalization remains a valid bridge from finite
words to actual positive-orbit dynamics, but its corridor-uniform closing hypothesis is, once the
corridor width is computed rather than assumed, essentially an $O(\log n)$ stopping-time
conjecture; and deep negative drift is self-financing, since the ordinary height it creates is
exactly the resource that permits the recovery valuations repaying it.

Neither EOC nor the Collatz conjecture is proved here. The surviving frontier is actual
positive-orbit arithmetic at three nested levels: qualitative exclusion of an infinite
zero-corridor tail, quantitative control of single-window lifetime, and control of separated
returns sufficient for full occupation.
\end{abstract}

\noindent\textbf{Non-claims.} This paper proves neither the Collatz conjecture nor any tier of EOC.
It does not prove an exponential realizer floor for arbitrary confined words. It does not promote
any computational finding to a theorem: every claim is labelled with its status
(\S\ref{sec:status}). Where a statement reorganizes or reclassifies earlier work rather than
proving something new, this is said explicitly.

\tableofcontents

\section{Introduction}\label{sec:intro}

The accelerated $3x+1$ map sends an odd positive integer to the next odd term of its Collatz
trajectory,
\[
T(m)=\frac{3m+1}{2^{a(m)}},\qquad a(m)=\nu_2(3m+1)\ge 1 .
\]
Writing $S_n$ for the accumulated $2$-adic valuation after $n$ accelerated steps and
$R_n=S_n-n\al$ with $\al=\log_2 3$ for its drift below the critical slope, the \emph{Effective
Occupation Conjecture} asserts that the time an orbit spends at or below any fixed drift level is
logarithmic in the seed. This single pointwise statement forces every orbit to be eventually
periodic and reduces Collatz to a finite cycle search (Theorem~\ref{thm:reduction}).

\subsection{Four axes, not three layers}

Revision~5 of this work was organized in three layers: EOC and its logical status; unconditional
drift-axis progress; and the residue-axis reformulation. That organization is no longer adequate,
for a specific reason: the residue axis and the confinement-lifetime question are not two problems
but two coordinates on one problem, and a fourth, genuinely qualitative question --- exclusion of
divergent orbits --- sits beside both. Revision~6 is therefore organized around four axes.

\begin{description}[leftmargin=2.2em,style=nextline]
\item[Occupation axis (\S\ref{sec:occupation})] The functional $\Occ_c(m_0)=\#\{n\ge 0:R_n\le c\}$,
  counting \emph{all} times in the corridor, including separated later returns. EOC lives here, and
  it is the strongest of the four.
\item[Residue axis (\S\ref{sec:residue})] The exact realizer congruence: which integers realize a
  given valuation word, and how small the least such integer, $\rmin(N,c)$, can be.
\item[Time axis (\S\ref{sec:time})] The single-window lifetime $L_c(m)$ and the drift exit time
  $\tau_c(m)$. Record chronology makes the residue and time axes \emph{exactly} dual.
\item[Divergence axis (\S\ref{sec:divergence})] What is unconditionally known about a hypothetical
  divergent orbit: sparsity of its value set, summability of reciprocals, and exclusion of
  logarithmic drift floors.
\end{description}

\noindent The programme map in Figure~\ref{fig:map} records the exact bridges. Every arrow in it is
a theorem, and \S\ref{sec:map} gives the reference for each.

\begin{figure}[t]
\centering
\begin{minipage}{0.92\textwidth}
\small
\[
\begin{array}{c}
\textbf{EOC}\ \ \Occ_c(m_0)=O(\log m_0)\\[2pt]
\big\downarrow\ \text{\scriptsize$\Occ_c\ge L_c$, strict (Rem.~\ref{rem:LvsO})}\\[2pt]
\text{single-window lifetime } L_c(m)
\qquad\Big|\qquad
\text{separated returns (uncontrolled)}\\[2pt]
\big\updownarrow\ \text{\scriptsize record chronology (Prop.~\ref{prop:inversion})}\\[2pt]
\text{least-realizer frontier } \rmin(N,c)\\[2pt]
\big\updownarrow\ \text{\scriptsize exact congruence (Prop.~\ref{prop:congruence})}\\[2pt]
\text{arithmetic placement of } r(D)
\end{array}
\]
\[
\begin{array}{c}
\text{hypothetical divergent orbit}
\ \xrightarrow{\ \text{Prop.~\ref{prop:recip}}\ } R_n\to-\infty
\ \xrightarrow{\ \text{last global max}\ } \text{infinite zero-corridor tail}\\[2pt]
\xrightarrow{\qquad} \tau_0(m^*)=\infty
\ \xrightarrow{\ \text{Thm.~\ref{thm:DEequiv}}\ }
\text{failure of universal drift exit}
\end{array}
\]
\end{minipage}
\caption{The programme map. Vertical double arrows are exact equivalences; the lower chain is the
qualitative divergence reduction. Every implication is referenced in \S\ref{sec:map}.}
\label{fig:map}
\end{figure}

\subsection{What Revision 7 changes}\label{sec:changelog}

Revision~7 is a \emph{consolidation of corrections}. Stated plainly: the verification record
improved, several claims were narrowed, and universal drift exit remains open. Nothing here proves
divergence exclusion or any tier of EOC.

\begin{enumerate}[leftmargin=1.6em]
\item \textbf{Verification improved.} Reciprocal summability (Prop.~\ref{prop:recip}) is no longer
  an unformalized input: it is now \emph{derived in Lean} from the windowed sparsity count
  (Rem.~\ref{rem:divledger}). The remaining external input in that chain is
  Theorem~\ref{thm:GTC} itself, and the open hypothesis $(\mathrm{DE})$ remains open; both are
  retained explicitly. Curry's note was also audited from the source text, and its two local
  defects --- a typographical slip in the light-case bound and a missing integrality step --- were
  identified and repaired; neither affects Proposition~\ref{prop:recip}.
\item \textbf{Separated returns, made precise} (Open Problem~\ref{op:F}). The entry lemma pins
  every re-entry threshold below $\al-1$ (Prop.~\ref{prop:entry}); a \emph{uniformly bounded}
  episode count is refuted constructively (Prop.~\ref{prop:episodes}); and the coarse per-episode
  accounting is shown lossy by a factor $\Omega(\log m_0)$ on the realized prefix
  (Rem.~\ref{rem:episodescope}). Finite-prefix statements are kept separate from total occupation
  throughout.
\item \textbf{The corridor-uniform hypothesis, clarified --- and a gap in it closed}
  (Rems.~\ref{rem:cyclegap}, \ref{rem:Uallscope}). Revision~6 read the hypothesis as giving an
  $O(\log m)$ bound on the total stopping time \emph{to $1$}. It does not: it bounds the
  \emph{injective} initial segment, so the orbit enters \emph{a} cycle within $O(\log m)$ steps,
  and arrival at $1$ needs cycle exclusion, which the hypothesis does not visibly supply. The EOC
  consequence is restated with that assumption explicit, and what survives without it is stated
  separately. Also: the analytic absorption inside Corollary~\ref{cor:Uall} is identified as the
  unformalized layer, and corridor-uniformity must not be conflated with a bound at one fixed
  corridor.
\item \textbf{Scope overstatements corrected.} The closed-route table no longer labels
  fixed-modulus \emph{terminality} as a formal result; and the cardinality barrier is stated for
  the \emph{known divergent-orbit} shadows, not for word-level criteria in general --- one such
  criterion is exact (Rem.~\ref{rem:zcrecounter}).
\item \textbf{Hierarchy strictness audited} (Rem.~\ref{rem:strictness}). ``No converse holds in
  general'' and ``Level~4 is strictly stronger than Level~3'' are withdrawn: a pointwise example
  with $\Occ_c>L_c$ shows the quantities differ, not that the universal bounds are inequivalent.
  For the corridor-uniform form the implication in fact reverses.
\item \textbf{Bibliography.} Four entries recovered their DOIs, and Ba\v{r}ina's entry --- which
  had merged two distinct papers under one title --- is split into
  \cite{Barina2021} and \cite{Barina2025}, each with its own metadata.
\end{enumerate}

Revision~6 is preserved unchanged; this revision supersedes it only where stated above.

\subsection{What Revision 6 added}

Revision~6 is primarily a \emph{synthesis}. Its main additions are:

\begin{enumerate}[leftmargin=1.6em]
\item A general transfer principle (\S\ref{sec:transfer}) turning any pointwise lifetime bound into
  a realizer floor, with the exponents worked out: logarithmic, polylogarithmic, polynomial and
  linear lifetime bounds give exponential, stretched-exponential, polynomial and linear frontiers
  respectively. EOC thereby becomes a \emph{graded} programme rather than an all-or-nothing one.
\item The qualitative time axis (\S\ref{sec:driftexit}): the drift exit time $\tau_0$, its exact
  integral characterisation $R_n>0\iff 3^n<2^{S_n}$, and the equivalence between universal drift
  exit and the exclusion of divergent orbits.
\item A boundaries section (\S\ref{sec:boundaries}) collecting the strongest no-go results: the
  finite-cylinder barrier, the coordinate-collapse ledger, and the word-shadow cardinality barrier.
\item A corrected reading of the launch-renormalization closing hypothesis (\S\ref{sec:launch}) and
  the self-financing obstruction to temporal impossibility arguments (\S\ref{sec:selffin}).
\item A closed-route map (\S\ref{sec:closed}) recording, for each mechanism tried, what it achieved
  and where it stops.
\end{enumerate}

We are explicit about provenance throughout. Several statements that read as new are in fact
consequences of identities already in Revision~5 --- the future-minimum theorem
(Prop.~\ref{prop:futuremin}) follows from the Eliahou--Rozier identity, and the positive cycle
drift (Lemma~\ref{lem:cycledrift}) is Revision~5's Lemma on terminal and cyclic drift. Where this
happens we say so and record only the new proof or the new role.

\section{Preliminaries}\label{sec:prelim}

Let $m_0$ be a positive odd integer, $m_{n+1}=T(m_n)$, $d_n=a(m_n)$, $S_n=\sum_{i<n}d_i$,
$\al=\log_2 3$, $R_n=S_n-n\al$, so $R_0=0$ and each step changes $R$ by $d_n-\al$.

\begin{remark}
If every accelerated orbit reaches $1$ then Collatz holds: every positive integer's odd part has an
accelerated orbit, and the classical trajectory interleaves accelerated iterates with halvings. We
work exclusively with $T$. \status{standard}
\end{remark}

For fixed $c>0$, $\Occ_c(m_0)=\#\{n\ge0:R_n\le c\}\in\mathbb{Z}_{\ge1}\cup\{\infty\}$. This is a
global sojourn count: it neither halts at first passage nor counts time far above the critical line.

\begin{lemma}[Eliahou--Rozier identity]\label{lem:ER}
For every odd $m_0\ge1$ and every $n\ge0$,
\[
m_n\,2^{R_n}=m_0\prod_{i<n}\Bigl(1+\tfrac{1}{3m_i}\Bigr),
\qquad\text{equivalently}\qquad
R_n=\log_2\frac{m_0}{m_n}+E_n,\quad
E_n=\sum_{i<n}\log_2\Bigl(1+\tfrac{1}{3m_i}\Bigr)\ge0 .
\]
\status{due to Eliahou \cite{Eliahou} and Rozier \cite{Rozier}; no novelty claimed}
\end{lemma}

\begin{remark}[Notation warning: two quantities called $E$]\label{rem:twoE}
Throughout, $E_n=\log_2 Q_n$ denotes the \emph{carry excess} of Lemma~\ref{lem:ER}. The
\emph{endpoint depth} of a word, elsewhere written $E(D)$, is a different quantity and is written
$\Edep(D)=S-\log_2 r(D)$ here (\S\ref{sec:gap}). The two are unrelated and were a genuine source of
confusion in earlier drafts.
\end{remark}

There is a second, purely integral form of the same information, used whenever real logarithms
would be inconvenient.

\begin{lemma}[Aggregate identity]\label{lem:aggregate}
With $C_n=\sum_{j<n}3^{\,n-1-j}2^{S_j}$ (so $C_0=0$, $C_{n+1}=3C_n+2^{S_n}>0$ for $n\ge0$),
\[
2^{S_n}m_n=3^n m_0+C_n \qquad\text{for all }n\ge0 .
\]
\status{elementary; formalized, \S\ref{sec:lean}}
\end{lemma}

\begin{proof}
Induction on $n$. The case $n=0$ is $m_0=m_0$. If $2^{S_n}m_n=3^nm_0+C_n$ then, multiplying
$2^{d_n}m_{n+1}=3m_n+1$ by $2^{S_n}$,
$2^{S_{n+1}}m_{n+1}=3\cdot 2^{S_n}m_n+2^{S_n}=3^{n+1}m_0+(3C_n+2^{S_n})=3^{n+1}m_0+C_{n+1}$.
\end{proof}

\begin{lemma}[Integral corridor]\label{lem:integral}
Since $S_n\in\mathbb{Z}$ and $n\al=\log_2(3^n)$,
\[
R_n\le 0\iff 2^{S_n}\le 3^n,\qquad R_n>0\iff 3^n<2^{S_n},\qquad R_n\le0\iff S_n\le\lfloor n\al\rfloor .
\]
\status{elementary}
\end{lemma}

\begin{lemma}[Wedge]\label{lem:wedge}
If $m_i\ge m_0$ for all $0\le i\le n$, then $R_i\le E_i\le i/(3m_0\ln2)$ for $0\le i\le n$. In
particular $R_i\le c$ for $0\le i\le\min\{n,\lfloor 3c\ln2\cdot m_0\rfloor\}$.
\status{Revision 5}
\end{lemma}

\section{The occupation axis: EOC and its logical status}\label{sec:occupation}

\subsection{Four tiers}

\begin{conjecture}[Existence EOC]\label{conj:exist}
For every fixed $c>0$ there exist constants $K_c,C_c$ such that every positive odd $m_0$ satisfies
$\Occ_c(m_0)\le K_c\log_2 m_0+C_c$.
\end{conjecture}

\begin{conjecture}[Effective EOC at $(c,K,C)$]\label{conj:eff}
For stated $c>0$ and stated numerical $K,C$, every positive odd $m_0$ satisfies
$\Occ_c(m_0)\le K\log_2 m_0+C$.
\end{conjecture}

\begin{conjecture}[Sharp-leading EOC]\label{conj:sharp}
For every fixed $c>0$, $\Occ_c(m_0)\le K\log_2 m_0+O_c(\log\log m_0)$ with $K=1/I(\al)$, where
$I(\mu)=\mu-\mu\log_2\mu+(\mu-1)\log_2(\mu-1)$.
\end{conjecture}

\begin{conjecture}[Strong sharp EOC]\label{conj:strong}
With $K=1/I(\al)$ as in Conjecture~\ref{conj:sharp}: for every fixed $c>0$ there is $C_c$ with
$\Occ_c(m_0)\le K\log_2 m_0+C_c$ for every positive odd $m_0$. (Revision~6 left $K$ unbound here.)
\end{conjecture}

Conjecture~\ref{conj:strong}$\Rightarrow$\ref{conj:sharp}$\Rightarrow$\ref{conj:exist}, and
Conjecture~\ref{conj:eff} at any constants implies Conjecture~\ref{conj:exist} at that $c$.
Crucially, Conjecture~\ref{conj:exist} does \emph{not} imply Conjecture~\ref{conj:eff} at any
stated constants: it asserts existence without a bound. All four tiers are pointwise statements
about every seed, not almost-every-seed statements.

\subsection{The Collatz reduction}

\begin{definition}\label{def:mstar}
For $c>0$ and constants $K,C$, let $m^*=m^*(K,C,c)$ be least such that
$3c\ln2\cdot m>K\log_2 m+C+1$ for all $m\ge m^*$; it exists and is computable by finitely many
exact comparisons.
\end{definition}

\begin{theorem}\label{thm:reduction}
Fix $c>0$ and suppose Conjecture~\ref{conj:exist} holds at $(c,K,C)$; put $m^*=m^*(K,C,c)$. Then
\begin{enumerate}[label=(\roman*),leftmargin=2em]
\item no accelerated orbit is injective; every orbit is eventually periodic;
\item every odd $m_0\ge m^*$ has an iterate $m_i<m_0$ with $i\le\lfloor 3c\ln2\cdot m_0\rfloor$;
\item if some seed's orbit does not reach $1$, there is a nontrivial accelerated cycle with minimum
  member below $m^*$;
\item if additionally no nontrivial cycle has minimum member below $m^*$, the Collatz conjecture
  holds.
\end{enumerate}
\status{Revision 5, retained unchanged}
\end{theorem}

\begin{corollary}[Explicit instance]\label{cor:explicit}
Assume Conjecture~\ref{conj:eff} at $c=1$, $K=13$, $C=100$. Then $m^*(13,100,1)=90$; every odd
$m<90$ reaches $1$ by direct inspection (longest: $m=73$, $42$ accelerated steps; $m=27$ and
$m=55$ take $41$); hence Collatz holds. \status{Revision 5; witness corrected in Revision 7}
\end{corollary}

\begin{remark}[Why effective constants are genuinely needed]
$m^*$ is computable from $(K,C,c)$, but Conjecture~\ref{conj:exist} supplies no bound on $(K,C)$
themselves. Minimal-counterexample reasoning does not repair this: Theorem~\ref{thm:reduction}(ii)
localizes a counterexample below $m^*$ but does not eliminate it. Since Ba\v{r}ina's
verification~\cite{Barina2025} settles all seeds below $2^{71}$, any nontrivial cycle's minimum member
exceeds $2^{71}$, and the wedge forces $\Occ_1(m_{\min})\gtrsim 3\ln2\cdot2^{71}\approx4.9\times10^{21}$
--- a value Conjecture~\ref{conj:exist} can accommodate at unspecified constants but which is
wildly larger than anything the record table of \S\ref{sec:comp} suggests. \status{Revision 5}
\end{remark}

\subsection{The entropy scale}\label{sec:entropy}

The natural population constant is the per-odd-step entropy deficit
\[
I_{\mathrm{Collatz}}\;=\;\al\bigl(1-H_2(1/\al)\bigr)\;=\;\frac{1-H_2(\rho)}{\rho},
\qquad \rho=\frac{1}{\al}=0.6309297\ldots,
\]
with $H_2$ the binary entropy. Numerically $I_{\mathrm{Collatz}}=0.0793186127748554\ldots$, and
$1/I(\al)=12.6074\ldots$ is the predicted asymptotic slope in Conjecture~\ref{conj:sharp}. The
parenthesization matters: the constant is $\bigl(1-H_2(\rho)\bigr)/\rho$, \emph{not}
$1-H_2(\rho)/\rho$. \status{arithmetic; a known typesetting hazard in earlier drafts}

This constant calibrates the population scale and supplies, by itself, no pointwise
anti-concentration. That distinction is one of the recurring themes of \S\ref{sec:boundaries}.

\subsection{EOC as an occupation-time shadow of Rozier's hypothesis}\label{sec:rozier}

\begin{theorem}\label{thm:rozier}
Assume Rozier's Lower Bound Hypothesis \cite[Hyp.~2.3]{Rozier}. Then for every fixed $c>0$,
$\Occ_c(m_0)\le \frac{1}{I(\al)}\log_2 m_0+O_c(\log\log m_0)$; that is, LBH implies
Conjecture~\ref{conj:sharp}, hence Conjecture~\ref{conj:exist}. \status{Revision 5}
\end{theorem}

The proof applies Rozier's own manoeuvre \cite[Thm.~4.1]{Rozier} at occupation times rather than at
total stopping time, via the entropy identity $\mu(1-H_2(1/\mu))=I(\mu)$. We regard it as a
corollary of his framework, recorded because the occupation functional had not previously been
isolated. It does not advance the provability of EOC --- LBH is itself open, and arguably harder ---
but positions EOC precisely as a weaker occupation-time shadow of LBH. Making the additive term
explicit would supply Conjecture~\ref{conj:eff} and, with it, Collatz via
Corollary~\ref{cor:explicit}; this is Open Problem~\ref{op:A}.

\section{The divergence axis}\label{sec:divergence}

This section is unconditional throughout; nothing here assumes EOC.

\begin{definition}
An orbit $(m_n)$ is \emph{injective} if $m_0,m_1,\dots$ are pairwise distinct.
\end{definition}

\begin{lemma}[Terminal and cyclic drift]\label{lem:cycledrift}
If an orbit reaches $1$, or enters a nontrivial positive cycle, then $R_N\to+\infty$. Consequently
every non-injective orbit has $R_N\to+\infty$. In particular, closure of a cycle of length $L$ with
total valuation $S$ forces $2^S>3^L$: strictly positive drift per period.
\status{Revision 5; the integral form $2^S>3^L$ is immediate from Lemma~\ref{lem:aggregate} and is
formalized, \S\ref{sec:lean}}
\end{lemma}

\subsection{An elementary $8/9$ bound and its terminality}

\begin{lemma}[Carry budget]\label{lem:carrybudget}
Every accelerated iterate $m_k$, $k\ge1$, satisfies $m_k\not\equiv0\pmod 3$. Hence for $k\ge1$ the
states lie in the density-$1/3$ set of integers coprime to $6$, and for every injective orbit
$E_N\le\frac19\log_2 N+O(1)$, so $2^{E_N}\ll N^{1/9}$. \status{Revision 5}
\end{lemma}

\begin{theorem}[Logarithmic-floor exclusion, $8/9$]\label{thm:89}
Let $m_0$ have an injective accelerated orbit. For every $B<8/9$, $R_k\ge-B\log_2 k$ fails for
infinitely many $k$. \status{Revision 5}
\end{theorem}

\begin{corollary}\label{cor:injocc}
For every injective orbit and fixed $c>0$, $\#\{0\le j<N:R_j>c\}\ll_{m_0,c}N^{1/9}$; hence
$\Occ_c(m_0)=\infty$ at density-one rate, and Conjecture~\ref{conj:exist} directly excludes
injective orbits. This is what Theorem~\ref{thm:reduction}(i) uses. \status{Revision 5}
\end{corollary}

\begin{remark}[Terminality of the mod-$3$ mechanism]\label{rem:terminal}
For a prime $p>3$, both $2$ and $3$ are units mod $p$, so $3m+1\equiv0\pmod p$ is solvable and no
annihilation analogous to the mod-$3$ case occurs. This one-step fixed-modulus packing mechanism
supplies no further density reduction at any other fixed modulus; pushing past $8/9$ requires a
genuinely different mechanism. \status{Revision 5}
\end{remark}

\subsection{Garcia--Tal/Curry sparsity}

\begin{theorem}[Garcia--Tal \cite{GarciaTal}; explicit form, Curry \cite{Curry}]\label{thm:GTC}
Let $\Occ$ be the value set of an infinite aperiodic orbit of the raw map. For every
$\beta>\beta_*$ there is $C_\beta$ with
$\#(\Occ\cap[a,a+X))\le C_\beta X^\beta\log(2X)$ uniformly in $a$, where
$\beta_*=\gamma_*\log_2 3\approx0.9653844$ and $\gamma_*\approx0.6090897$ is the unique root of
$H(\gamma)=\gamma\log_2 3$. \status{external, cited}
\end{theorem}

\begin{proposition}[Reciprocal summability]\label{prop:recip}
Every divergent (equivalently, injective) orbit satisfies $\sum_{x\in\Occ}1/x<\infty$, hence
$\sum_{n\ge0}1/m_n<\infty$ for its accelerated iterates. Consequently $Q_n\uparrow Q_\infty<\infty$,
$E_n\uparrow E_\infty<\infty$, $m_n\to\infty$, and $R_n\to-\infty$. \status{Revision 5, from
Theorem~\ref{thm:GTC}}
\end{proposition}

\begin{theorem}[Occupation bound and the $1/\beta_*$ threshold]\label{thm:betastar}
Fix $\beta\in(\beta_*,1)$. There are constants $c_1,c_2$ depending on $\beta$ and on the orbit with
$\#\{n\ge0:R_n\ge-G\}\le c_1 2^{\beta G}(G+c_2)$ for all $G\ge1$. Consequently no divergent orbit
satisfies $R_n\ge-B\log_2 n+O(1)$ eventually, for any $B<1/\beta_*\approx1.0358567$. The endpoint
$B=1/\beta_*$ is not excluded. \status{Revision 5}
\end{theorem}

\subsection{Curry's second, qualitative role}\label{sec:curryqual}

Revision~5 treated Theorem~\ref{thm:GTC} purely as quantitative divergence sparsity. It has a
second role, which organizes the whole time axis.

\begin{proposition}[Last-maximum reduction]\label{prop:lastmax}
Let $m_0$ have a divergent orbit. By Proposition~\ref{prop:recip}, $R_n\to-\infty$, so $R$ attains
a global maximum at finitely many indices; let $n_0$ be the \emph{last}. Put $m^*=m_{n_0}$, an
actual positive odd integer, and $R'_k=R_{n_0+k}-R_{n_0}$ its relative drift. Then $R'_k<0$ for
every $k\ge1$, hence $S'_k\le\lfloor k\al\rfloor$ for every $k\ge1$: the restarted orbit has an
infinite zero-corridor tail. \status{companion work. Formalized as
\texttt{exists\_zero\_confined\_seed\_tendsto}, but under the hypothesis
\texttt{ReciprocalSummable}\,---\,i.e. conditional on Theorem~\ref{thm:GTC}, which is external and
unformalized. It is \emph{not} formalized from ``divergent orbit'' directly}
\end{proposition}

This reduction is not new to Revision~6 --- it is the zero-corridor tail theorem of the Curry
foundation notes, where it is proved and formalized, and where it is also recorded that the step
requires no irrationality of $\al$. Its role here is to make the divergence axis and the time axis
meet.

\begin{remark}[Scope]
Nothing in this section constrains hypothetical nontrivial cycles, proves any tier of EOC, or
touches the occupation of convergent orbits. \status{Revision 5}
\end{remark}

\section{The residue axis}\label{sec:residue}

\begin{definition}
A valuation word $D=(d_0,\dots,d_{N-1})$ has $d_i\ge1$; $s_0=0$, $s_j=\sum_{i<j}d_i$,
$S_N(D)=s_N$, and carry sum $C_N(D)=\sum_{j=0}^{N-1}3^{\,N-1-j}2^{s_j}$. An odd $m_0$
\emph{realizes} $D$ if $a(m_i)=d_i$ for all $i<N$.
\end{definition}

\begin{proposition}[Exact realizer congruence]\label{prop:congruence}
The exact realizers of $D$ form a single residue class modulo $2^{S_N(D)+1}$, not $2^{S_N(D)}$.
Writing $r(D)$ for the least positive representative,
\[
r(D)\equiv\bigl(2^{S_N(D)}-C_N(D)\cdot3^{-N}\bigr)\pmod{2^{S_N(D)+1}} .
\]
\status{Revision 5}
\end{proposition}

\begin{definition}[Coarse anchor versus exact lift]
For $S\le S_N(D)$ the \emph{coarse anchor} at depth $S$ is
$a_S(D):=\bigl(-C_N(D)3^{-N}\bigr)\bmod 2^S$. Proposition~\ref{prop:congruence} shows $r(D)$
reduces mod $2^{S_N}$ to $a_{S_N}(D)$ and is one of $\{a_{S_N}(D),a_{S_N}(D)+2^{S_N}\}$, terminal
parity selecting which. Unconditionally $r(D)\ge a_{S_N}(D)$; this inequality, not the equality, is
what every floor argument below uses. \status{Revision 5}
\end{definition}

\subsection{Proved sectors}

\begin{theorem}[Periodic realizer floor]\label{thm:periodic}
For a fixed block $X$ put $\xi_X=-C(X)/\delta_X\in\Zt$ with $\delta_X=3^{q_X}-2^{S_X}$ odd. If
$\xi_X\notin\Zpos$, there is a finite $B_X$ with
$r(X^n)\ge 2^{nS_X-B_X}$ for all large $n$, hence $\log_2 r(X^n)/(nS_X)\to1$. \status{Revision 5;
full proof in companion note \cite{DJperiodic}}
\end{theorem}

\begin{theorem}[Eventually periodic words]\label{thm:evperiodic}
For a fixed prefix $Y$ the same mechanism applies to $YX^n$ with transformed anchor
$\xi_{X,Y}$; outside the same exceptional case, $r(YX^n)\ge2^{S_Y+nS_X-B_{X,Y}}$. \status{Revision 5}
\end{theorem}

\begin{proposition}[Bulk-defect anchor stability]\label{prop:bulkdefect}
For $D=X^aEX^b$ (fixed $X$, arbitrary finite defect $E$, arbitrary $b$), the coarse anchor satisfies
$a_{aS_X}(D)=\xi_X\bmod2^{aS_X}$ exactly, independent of $E$ and $b$. Hence, outside the
exceptional case, $r(D)\ge2^{aS_X-B_X}$. \status{Revision 5}
\end{proposition}

\begin{remark}[Arbitrary irregular words: open]\label{rem:moving}
For a general confined word the coarse anchor $\xi_N(D)=-C_N(D)3^{-N}$ changes with $N$: there is
no single fixed rational whose periodicity can be invoked. This is the moving-anchor sector, and it
is where actual Collatz orbits live. \status{Revision 5}
\end{remark}

\subsection{What the height invariant actually separates}\label{sec:height}

Revision~5 framed the split as periodic versus irregular. Later audits sharpened it. Every realizer
condition can be written
\[
2^{S_N}\ \big|\ A+\chi\Lambda
\]
for integers $A,\Lambda$ \emph{generated by the word}: for a general word $(A,\Lambda)=(C_N,3^N)$,
while for a purely periodic word of period $p$ the carry collapses to
$C_{Np}\cdot\Delta=c_0(3^{Np}-2^{N\sigma})$ with $c_0,\Delta$ independent of $N$. The only general
tool converting divisibility into a size statement is that a nonzero multiple of $2^m$ has absolute
value at least $2^m$, which gives a \emph{dichotomy}: either $A+\chi\Lambda=0$, pinning $\chi$ to
the single rational $-A/\Lambda$, or
\[
|\chi|\ \ge\ 2^{S_N}/H-1,\qquad H=\max(|A|,|\Lambda|).
\]
Writing $\rho_N=\log_2 H/S_N$, the second branch is useful exactly when $\rho_N$ is bounded away
from $1$.

\begin{observation}[Height amortization, not complexity]\label{obs:height}
Periodic, eventually periodic and one-defect words have $\rho_N\to0$, at rate $1/(\text{number of
repetitions})$: a fixed height is amortized over unboundedly many periods. Arbitrary confined words
supply only $(C_N,3^N)$, with $\log_2 H=\al N\ge S_N$, so $\rho_N\ge1$ and the second branch is
vacuous. The separating invariant is therefore height amortization, not symbolic complexity: the
Beatty/Sturmian zero-corridor word, of minimal aperiodic subword complexity $p(n)=n+1$, already has
$\rho_N>1$ at every depth examined, while its periodic continued-fraction approximants keep
$\rho_N\approx0.26$. \status{structural interpretation, supported by exact identities and by
computation; not a theorem about all words}
\end{observation}

\subsection{Endpoint depth and the frontier defect}\label{sec:gap}

Write $\Edep(D)=S-\log_2 r(D)$ for the endpoint depth (Remark~\ref{rem:twoE}) and
$g(D)=S-\lfloor\log_2 r(D)\rfloor$ for the discrete terminal gap. Then
\[
\Edep(D)\ \le\ g(D)\ <\ \Edep(D)+1,\qquad\text{i.e. } g=\lceil \Edep\rceil ,
\]
the two agreeing only when $r(D)$ is a power of $2$.
\status{proved; the inequalities are formalized
(\texttt{discreteGap\_bracket}); the equality clause is not}

\begin{remark}[Correction to Revision~5's Open Problem E]\label{rem:OPE}
Revision~5 asked whether the frontier-defect target
$\Edep(D)\le H_2(\rho_c)S+\Phi(N)$, $\Phi=o(N)$, is genuinely easier than an exponential floor, and
recorded that question as open. It is not easier in any useful sense: by the bracket above the
statement is equivalent, up to an additive $1$, to $g(D)\le H_2(\rho_c)S+\Phi(N)$, and under
confinement $S\le\lfloor\al N\rfloor$ this is an exponential floor with
$\varepsilon=I_{\mathrm{Collatz}}+o(1)$. Open Problem~E of Revision~5 is therefore the
sharp-leading form of Open Problem~C, not an independent weaker route, and is restated as such in
\S\ref{sec:open}. \status{correction to Revision 5}
\end{remark}

\section{The time axis}\label{sec:time}

\subsection{Three temporal scales}

\begin{definition}\label{def:times}
Fix $c\ge0$ and an odd seed $m$.
\begin{align*}
L_c(m)&=\max\{N\ge0: R_j(m)\le c\ \text{for all }0\le j\le N\} &&\text{(single-window lifetime)},\\
\Occ_c(m)&=\#\{n\ge0:R_n(m)\le c\} &&\text{(total occupation)},\\
\tau_c(m)&=\inf\{n\ge1:R_n(m)>c\} &&\text{(drift exit time)} .
\end{align*}
\end{definition}

These answer different questions. $L_c$ measures the initial confined episode; $\Occ_c$ sums all
episodes, including separated later returns; $\tau_c$ asks only whether the corridor is ever left.
Always $L_c(m)\le\Occ_c(m)$ and $\tau_c(m)=L_c(m)+1$ when finite.

\begin{remark}[Single window does not control occupation]\label{rem:LvsO}
The inequality $\Occ_c\ge L_c$ is strict in general, and dramatically so: $m_0=285175$ has
$L_1=14$ but $\Occ_1=97$. So the two quantities are genuinely different functionals, and a
single-window theory does not \emph{directly} deliver EOC.

Revision~6 drew the stronger conclusion that ``a complete single-window theory --- even an optimal
one --- does not yield EOC''. That does not follow from a pointwise example, and is withdrawn; see
Remark~\ref{rem:strictness}. For the corridor-uniform form of the hypothesis it is in fact false.
\status{Revision 5 example; conclusion corrected in Revision 7}
\end{remark}

\subsection{Residue--time duality}\label{sec:duality}

\begin{definition}
Fix $c>0$. A word $D$ of length $N$ is \emph{$c$-confined} if $R_j(D)\le c$ for $0\le j\le N$. For
the $c$-confined family at depth $N$, $\rmin(N,c):=\min_D r(D)$.
\end{definition}

\begin{proposition}[Record-chronology inversion]\label{prop:inversion}
$\rmin(N,c)=\min\{m\ \text{odd}: L_c(m)\ge N\}$. \status{Revision 5, Proposition 5.6}
\end{proposition}

\begin{proposition}[Single-window equivalence]\label{prop:window}
Fix $c>0$ and $\varepsilon\in(0,2]$. The following are equivalent:
\begin{enumerate}[label=(\roman*),leftmargin=2em]
\item $r(D)\ge2^{\varepsilon N}$ for every $c$-confined $D$ of length $N$ with $r(D)\ge3$;
\item $L_c(m_0)\le\varepsilon^{-1}\log_2 m_0$ for every odd $m_0\ge3$.
\end{enumerate}
\status{Revision 5, Proposition 5.7}
\end{proposition}

Revision~5 presented Propositions~\ref{prop:inversion} and~\ref{prop:window} as technical
observations inside the residue section. They deserve a more prominent role: together they say that
\emph{the residue and time axes are the same axis in two coordinates}. Every statement about how
small a realizer can be is a statement about how long an orbit can stay confined, and conversely.
Revision~6 takes this as an organizing principle rather than a lemma.

\subsection{The general transfer principle}\label{sec:transfer}

Proposition~\ref{prop:window} matches one specific pair of scales. The inversion is general.

\begin{proposition}[Lifetime--frontier transfer]\label{prop:transfer}
Let $F:\mathbb{Z}_{>0}\to\mathbb{Z}_{>0}$ be nondecreasing and suppose $L_c(m)\le F(m)$ for every
odd $m$. Then
\[
\rmin(N,c)\ \ge\ \min\{m: F(m)\ge N\}\;=\;F^{-1}(N).
\]
\status{proved here; the finite form is formalized, \S\ref{sec:lean}}
\end{proposition}

\begin{proof}
By Proposition~\ref{prop:inversion}, $\rmin(N,c)$ is the least odd $m$ with $L_c(m)\ge N$. If
$m<F^{-1}(N)$ then $F(m)<N$ by monotonicity, so $L_c(m)\le F(m)<N$; hence no such $m$ qualifies.
\end{proof}

Writing the hypothesis in terms of $x=\log_2 m$ gives the scales of interest.

\begin{center}
\begin{tabular}{@{}lll@{}}
\toprule
Lifetime hypothesis & Frontier obtained & Character \\
\midrule
$L_c(m)\le K\log_2 m$ & $\log_2\rmin(N,c)\ge N/K$ & exponential (Open Problem~C) \\
$L_c(m)\le K(\log_2 m)^p$, $p>1$ & $\log_2\rmin(N,c)\ge(N/K)^{1/p}$ & stretched exponential \\
$L_c(m)\le K m^{\theta}$, $\theta>0$ & $\rmin(N,c)\ge(N/K)^{1/\theta}$ & polynomial \\
$L_c(m)\le K m$ & $\rmin(N,c)\ge N/K$ & linear \\
\bottomrule
\end{tabular}
\end{center}

\begin{remark}[A graded programme]\label{rem:graded}
This table is the most constructive consequence of the recent audits. EOC's residue axis has
usually been posed all-or-nothing --- prove an exponential floor or prove nothing. Transfer shows
otherwise: \emph{any} unconditional pointwise lifetime bound, however weak, yields a corresponding
unconditional frontier bound. A bound as weak as $L_c(m)\le Km$ would already give a linear
frontier, and none is currently known. We record this as Open Problem~\ref{op:intermediate}.
\status{proved; the assessment of what is currently known is in \S\ref{sec:nounconditional}}
\end{remark}

\begin{remark}[No unconditional lifetime bound can be easy]\label{sec:nounconditional}
No unconditional bound $L_0(m)\le F(m)$ is currently known, and none can be proved without
excluding divergent orbits: applied to the restarted seed $m^*$ of
Proposition~\ref{prop:lastmax}, which has $L_0(m^*)=\infty$, any such bound is immediately
contradicted. The transfer table therefore grades a programme whose every entry is open.
\status{proved}
\end{remark}

\subsection{Drift exit and the future-minimum theorem}\label{sec:driftexit}

\begin{proposition}[Future minimum]\label{prop:futuremin}
For $n\ge1$: if $R_n\le0$ then $m_n>m_0$. Hence an orbit with $R_n\le0$ for all $n\ge1$ has its
seed as a \emph{strict minimum of its entire future}. Contrapositively, for $n\ge1$, $m_n\le m_0$
already forces $R_n>0$.

The restriction $n\ge1$ is necessary, not cosmetic: at $n=0$ we have $R_0=0\le0$ while
$m_0>m_0$ fails. Both proofs below use it --- the first through $E_n>0$, the second through
$C_n>0$, and $E_0=C_0=0$. \status{Revision 5; hypothesis $n\ge1$ made explicit in Revision 7}
\end{proposition}

\begin{proof}
Two proofs. Analytically, Lemma~\ref{lem:ER} gives $R_n=\log_2(m_0/m_n)+E_n$ with $E_n>0$ for
$n\ge1$, so $m_n\le m_0$ implies $R_n\ge E_n>0$. Integrally, Lemma~\ref{lem:aggregate} and
$C_n>0$ give $2^{S_n}m_n=3^nm_0+C_n>3^nm_0\ge2^{S_n}m_0$ when $2^{S_n}\le3^n$, so $m_n>m_0$.
\end{proof}

\begin{remark}[Provenance]
This is \emph{not} a new theorem. The analytic proof is an immediate consequence of the
Eliahou--Rozier identity already in Revision~5 (Lemma~\ref{lem:ER}), and is close in content to
Revision~5's wedge (Lemma~\ref{lem:wedge}). What is new is only the second proof, which uses no
logarithms and is therefore convenient to formalize, and the role the statement plays below.
\status{reclassification}
\end{remark}

\begin{proposition}[Drift exit is weaker than ordinary descent]\label{prop:weaker}
Let $\sigma(m)=\inf\{n\ge1:m_n<m\}$ be the (Terras) stopping time. Then $m_n<m_0\iff R_n>E_n$,
whence $\tau_0(m)\le\sigma(m)$, and $\sigma(m)<\infty\Rightarrow\tau_0(m)<\infty$. The inequality
is strict at cycle minima: for $m=1$, $R_n=E_n=(2-\al)n$ exactly, so $\tau_0(1)=1$ while
$\sigma(1)=\infty$. \status{proved}
\end{proposition}

\begin{remark}[How much weaker, quantitatively]\label{rem:howweak}
Off cycle minima the gap is small. At $n=\tau_0$, strictness needs $R_n\le E_n$; under confinement
for $j<n$, Proposition~\ref{prop:futuremin} gives $m_j>m_0$, so
$E_n\le(3\ln2)^{-1}\sum_{j<n}1/m_j\le 0.4809\,n/m_0$, while integrality of $S_n$ with $R_n>0$
forces $R_n\ge1-\{\al n\}$. Strictness at depth $n$ therefore requires
\[
\rmin(n,0)\ \le\ m_0\ \le\ 0.4809\,n/(1-\{\al n\}),
\]
whose right side is polynomial in $n$ because $\al$ has finite irrationality measure. So
\emph{any} superpolynomial realizer floor forces $\tau_0=\sigma$ away from cycle minima.
An exhaustive search of every odd $m$ below this size bound at every depth $n<200$ where the bound
is not already violated found no strict example. \status{the inequality is proved; the search is a
computational diagnostic}
\end{remark}

\subsection{Universal drift exit and divergence}\label{sec:DE}

\begin{definition}[Universal drift exit]
\[
(\mathrm{DE})\qquad\text{for every odd }m\ge1\text{ there is }n\ge1\text{ with } 3^n<2^{S_n(m)} .
\]
Equivalently, by Lemma~\ref{lem:integral}, $\tau_0(m)<\infty$ for every $m$. The phrase ``drift
exit time'' is terminology used here; we found no standard name for this quantity in the
literature (\S\ref{sec:prior}).
\end{definition}

\begin{theorem}[Divergence equivalence]\label{thm:DEequiv}
$(\mathrm{DE})$ holds if and only if there is no divergent positive accelerated orbit.
\end{theorem}

\begin{proof}
($\Rightarrow$, contrapositive.) If a divergent orbit exists, Proposition~\ref{prop:lastmax}
produces an actual positive odd integer $m^*$ with $R'_k<0$ for every $k\ge1$, so
$\tau_0(m^*)=\infty$ and $(\mathrm{DE})$ fails.

($\Leftarrow$, contrapositive.) Suppose $(\mathrm{DE})$ fails at $m$, so $2^{S_n}\le3^n$ for all
$n\ge1$. The orbit cannot repeat a state: a repeat makes it eventually periodic, and
Lemma~\ref{lem:cycledrift} then gives $R_N\to+\infty$, contradicting $R_n\le0$. So the orbit is
injective, hence divergent.
\end{proof}

\begin{corollary}[Cycle compatibility]\label{cor:cycles}
By Lemma~\ref{lem:cycledrift} every point of a nontrivial positive cycle, and every transient point
entering one, has finite $\tau_0$; so such a cycle is consistent with the defining condition of
$(\mathrm{DE})$ holding at every seed. Hence the drift-exit argument supplies no cycle exclusion,
and $(\mathrm{DE})$ alone does not reach the Collatz conjecture, which also excludes cycles.
\end{corollary}

\begin{remark}[``Strictly weaker'' is itself open]\label{rem:strictlyweaker}
Revisions~5--7 up to this point wrote that $(\mathrm{DE})$ is \emph{strictly weaker} than Collatz.
That is \textbf{not established}, and the argument just given does not establish it. Since
$(\mathrm{DE})$ is equivalent to the absence of divergent orbits (Thm.~\ref{thm:DEequiv}) and
Collatz is the conjunction of that with the absence of nontrivial cycles, Collatz
$\Rightarrow(\mathrm{DE})$ always, while $(\mathrm{DE})\Rightarrow$ Collatz holds \emph{if and only
if} no nontrivial cycle exists. So $(\mathrm{DE})$ is strictly weaker precisely when a nontrivial
cycle exists --- an open question whose expected answer is \emph{no}, in which case the two are
equivalent.

What is established, precisely, is this list and no more:

\begin{enumerate}[label=(\roman*),leftmargin=2em]
\item Collatz $\Rightarrow(\mathrm{DE})$;
\item $(\mathrm{DE})$ excludes divergence (Thm.~\ref{thm:DEequiv});
\item $(\mathrm{DE})$ together with the exclusion of nontrivial positive cycles $\Rightarrow$
  Collatz;
\item every point of a cycle has finite $\tau_0$ (Lem.~\ref{lem:cycledrift}), so the drift-exit
  argument itself supplies no cycle exclusion.
\end{enumerate}

Accordingly: \textbf{no implication from $(\mathrm{DE})$ to cycle-freeness is established here.}
We claim neither that $(\mathrm{DE})$ is strictly weaker than Collatz nor that $(\mathrm{DE})$
fails to entail cycle-freeness --- the latter is \emph{also} unproved, and would be false if no
nontrivial cycle exists, which is the expected case. Item~(iv) is a statement about what this
argument delivers, not about what is true. This is the same error of form corrected in
Remark~\ref{rem:strictness}: a non-implication of \emph{mechanism} is neither strictness nor
non-entailment of \emph{content}.
\status{Revision 7, correction of scope}
\end{remark}

\begin{remark}[Provenance, emphatically]\label{rem:DEprov}
Neither direction is new mathematics. The forward direction is the zero-corridor tail theorem of
the companion Curry foundation, already proved and formalized there; the reverse direction is
Revision~5's Lemma~\ref{lem:cycledrift}. Moreover the \emph{statement} $(\mathrm{DE})$ is, in
negated form, exactly the frontier already posed in that companion work: ``determine whether an
infinite zero-confined valuation word can be realized by a positive ordinary integer.'' Revision~6
contributes the reclassification --- recognising this as the qualitative endpoint of the
\emph{time} axis rather than an isolated realizability question --- together with the integral form
and its exact position in the hierarchy of \S\ref{sec:hierarchy}. We do not coin a new acronym for
it; in particular it is distinct from the repository's ZCRE statement, which is an equivalence
about bounded prefix realizers. \status{reclassification, not a new theorem}
\end{remark}

\section{Launch renormalization: the bridge to actual orbits}\label{sec:launch}

Finite word and residue analysis becomes powerful exactly when it produces an \emph{actual}
positive-orbit segment. Launch renormalization does this. We summarize only what the programme map
needs; the full treatment is the companion note \cite{DJlaunch}.

\begin{theorem}[Launch-state collapse, {\cite{DJlaunch}}]\label{thm:launch}
Let $D$ be a word and $j$ any index with $2^{S_j}>r(D)$. Then
$m^*_j:=(3^j r(D)+C_j(D))/2^{S_j}$ is a positive integer, in fact the genuine $j$-th accelerated
iterate of $r(D)$, and the letters $(d_{j+1},\dots,d_N)$ are exactly the accelerated valuation word
of $m^*_j$. \status{companion work}
\end{theorem}

Call $D$ of length $N$ \emph{$\varepsilon$-dangerous} if it is $c$-confined with
$\log_2 r(D)<\varepsilon N$. For such a word the pinning index $t_1$ and the launch data satisfy
\[
t_1\le\varepsilon N+O(1),\qquad
L\ \ge\ (1-\varepsilon)N-O(1),\qquad
\log_2 m^*_{t_1}\le \al\varepsilon N+O(\log N),
\]
with inherited \emph{relative} corridor $A\le c+(\al-1)\varepsilon N+O(1)$. The leading coefficient
in the launch-size bound is $\al\varepsilon$, not $\varepsilon$. The residual irregular class is
renormalization-stable: the forced suffix of a dangerous word in that class is again a small
integer with a long, near-critically confined injective orbit segment, so finite structural
enumeration cannot close the programme \cite{DJlaunch}.

\subsection{Reclassifying the corridor-uniform hypothesis}\label{sec:Uall}

The companion note closes the route conditionally on a corridor-uniform lifetime hypothesis: for
some $K<\infty$, every injective initial segment of $m$ with $R_i\le A$ for $0\le i\le L$ has
$L\le K(\log_2 m+A)$. This reads as a mild statement about corridors. It is not.

\begin{proposition}[The corridor is automatic]\label{prop:corridorauto}
For every positive orbit and every $n$, $2^{S_n}\le3^nm_0+C_n$, i.e.
$R_n\le\log_2 m_0+E_n$; and equality holds as soon as the orbit reaches $1$. Hence the least valid
corridor width for a full orbit is exactly $\log_2 m_0+E$, not merely bounded by it.
\status{proved; formalized, \S\ref{sec:lean}}
\end{proposition}

\begin{proof}
Immediate from Lemma~\ref{lem:aggregate} and $m_n\ge1$, with equality when $m_n=1$.
\end{proof}

\begin{corollary}[Reclassification]\label{cor:Uall}
Instantiating the corridor-uniform hypothesis at its own automatic corridor gives
$L\le K(2\log_2 m+E)$. For an injective segment the iterates are distinct odd integers, so
$E_L\le(3\ln2)^{-1}\sum_{j<L}1/m_j\le0.2404\ln L+O(1)$, which grows slower than $L$ and therefore
bounds $L$. Consequently the hypothesis implies that no accelerated orbit is injective forever ---
i.e. it implies the absence of divergent orbits --- and, more precisely, that \emph{every} orbit
repeats a value, and so enters a cycle, within $O(\log m)$ steps. Conversely a bound
$L\le K'\log_2 m$ returns the hypothesis for every $A\ge0$. \status{proved}
\end{corollary}

\begin{remark}[What the injective-segment bound does \emph{not} give]\label{rem:cyclegap}
The quantity bounded is the length of an \emph{injective} initial segment, which is the pre-period
plus the period: the orbit repeats a value within $O(\log m)$ steps and is thereafter periodic. It
does \emph{not} follow that the orbit reaches $1$: a hypothetical nontrivial positive cycle is
entirely consistent with the bound, and the corridor-uniform hypothesis does not visibly exclude
one. (Applied to a cycle point it bounds that cycle's \emph{period} by $O(\log m^*)$, which is a
strong constraint but not a contradiction.)

Revision~6 wrote ``and, with that, an $O(\log m)$ bound on accelerated total stopping time'', where
the total stopping time is the time to reach $1$. That step is \textbf{withdrawn}: it silently
assumed cycle exclusion. Everything downstream that used arrival at $1$ is restated in
Remark~\ref{rem:Uallscope} with the assumption made explicit. \status{Revision 7, correction}
\end{remark}

\begin{proposition}[Corridor times are large-value times]\label{prop:corridorvalue}
For every positive odd seed and every $n$, if $R_n\le c$ then $m_n\ge m_02^{-c}$.

\emph{Proof.} By Lemma~\ref{lem:aggregate}, $2^{S_n}m_n=3^nm_0+C_n\ge3^nm_0$, while $R_n\le c$
reads $2^{S_n}\le2^{c}3^n$. Hence $3^nm_0\le2^{c}3^nm_n$. \qed

No cycle hypothesis is used, and none is available: this is the only statement of its kind we can
make unconditionally. \status{Revision 7; formalized, \S\ref{sec:lean}}
\end{proposition}

\begin{proposition}[Occupation after the orbit becomes periodic]\label{prop:occstop}
Suppose the orbit of $m_0$ is eventually periodic, entering a cycle at time $T$ with value set
$\mathcal{C}$. Then
\[
\Occ_c(m_0)\ \le\ T+\#\{n\ge T:\ m_n\ge m_02^{-c}\},
\]
and the second term is $0$ as soon as $\max\mathcal{C}<m_02^{-c}$.

In particular, if the orbit reaches $1$ at time $n^*\ge1$ and $m_0>2^{c}$, then
$\Occ_c(m_0)\le n^*$.

\emph{Proof.} Immediate from Proposition~\ref{prop:corridorvalue}; for the last clause,
$\mathcal{C}=\{1\}$ and $1<m_02^{-c}$ exactly when $m_0>2^{c}$. \qed

\begin{remark*}
The hypotheses $n^*\ge1$ and $m_0>2^{c}$ (rather than $m_0\ge2^{c}$) are both needed, and the
boundary case shows why: for $m_0=1$, $c=0$ we have $n^*=0$, yet $R_0=0\le c$, so time $0$ lies in
the corridor and $\Occ_0(1)=1>n^*$. Proposition~\ref{prop:corridorvalue} diagnoses this exactly ---
at $m_0=1,c=0$ the criterion $m_n\ge m_02^{-c}$ is satisfied by the value $1$ itself. Revision~6's
formulation omitted these hypotheses.

What the Lean statement carries is $1\le n^*$ together with the \emph{non-strict} $2^{c}\le m_0$;
that pair suffices, and it is not the same pair as stated above. Revision~7 first claimed the Lean
hypothesis was $m_0>2^{c}$; it is not.
\end{remark*}
\status{Revision 7. The reaches-$1$ corollary is formalized
(\texttt{occupation\_le\_of\_reaches\_one}); the general eventually-periodic
display is \emph{not} formalized}
\end{proposition}

\begin{remark}[What Corollary~\ref{cor:Uall} does and does not settle]\label{rem:Uallscope}
Three points, kept apart.

\emph{(i) The hypothesis implies EOC --- given cycle exclusion.} Corollary~\ref{cor:Uall} bounds
the injective initial segment, i.e. the time $T$ at which the orbit first repeats a value, by
$O(\log m_0)$. Proposition~\ref{prop:occstop} then gives
\[
\Occ_c(m_0)\ \le\ T+\#\{n\ge T:\ m_n\ge m_02^{-c}\} ,
\]
whose second term vanishes once every value of the cycle entered is below $m_02^{-c}$. If the only
positive cycle is $\{1\}$ --- that is, \emph{given cycle exclusion} --- this holds for every
$m_0>2^{c}$, and the corridor-uniform hypothesis implies Existence EOC at every $c$. It then
relocates the difficulty rather than reducing it, since it also implies the Collatz conjecture.

\emph{Without} cycle exclusion the conclusion is weaker and should be stated as such: occupation is
bounded by $O(\log m_0)$ plus the corridor time spent inside whatever cycle the orbit enters. That
residual is finite for each fixed cycle --- a cycle has strictly positive period drift
(Lem.~\ref{lem:cycledrift}), so its drift exceeds $c$ after finitely many periods --- but we have no
bound on it that is uniform over hypothetical nontrivial cycles, whose period and values are
unconstrained above. Supplying such a uniform bound, or excluding the cycles, is a genuine
additional obligation. Revision~6 stated (i) without it. \status{Revision 7, correction}

\emph{(ii) One layer is not formalized.} Proposition~\ref{prop:occstop} and the aggregate identity
it rests on are machine-checked. The passage from the corridor-uniform hypothesis to the
stopping-time bound --- Corollary~\ref{cor:Uall} --- is not: it absorbs the term
$E_L=O(\log L)$ into a bound on $L$, an elementary but genuinely real-analytic step. The chain is
therefore a formalized bridge sitting on an unformalized analytic layer, and should be described
that way.

\emph{(iii) Do not conflate corridor-uniformity with a fixed corridor.} The hypothesis quantifies
over \emph{every} corridor width $A$ simultaneously; it is that uniformity which permits
instantiation at the automatic corridor $A=\log_2m_0+E$ and hence the collapse in (i). A bound on
$L_c$ at one fixed $c$ carries none of this: it is implied by the corridor-uniform form, and the
converse is not known. (We avoid ``strictly weaker'' here for the reason given in
Rem.~\ref{rem:strictlyweaker}.) Remarks elsewhere
in this paper about ``single-window'' bounds refer to the fixed-corridor form.
\status{Revision 7}
\end{remark}

So the corridor-uniform hypothesis is \emph{noncircular} with respect to Open Problem~C --- it is
not a repackaging of the realizer floor --- but it is a substantially stronger \emph{quantitative}
hypothesis than is required merely to exclude divergence. In particular, once the automatic
corridor is inserted it is essentially an $O(\log m)$-scale pointwise stopping/lifetime hypothesis,
and it should not be presented as a plausibly easier closing hypothesis. The
companion note's weakened variant, which exempts segments majority-covered by a bounded-period
terminal repetition, does not change this: over all odd $m<60001$ the exemption applies to none of
the $29\,999$ orbits tested. \status{computational diagnostic only}

\section{Self-financing recovery}\label{sec:selffin}

A natural strategy for the time axis is to argue that a deep negative drift excursion cannot be
repaid, because the recovery would demand an impossibly large valuation. This fails at its source.

\begin{proposition}[Self-financing]\label{prop:selffin}
A single step of valuation $d$ from state $m_n$ achieves drift exit as soon as
$2^{d}\,m_0>3\,m_n$. The available valuation is capped by $2^{d_n}\le 3m_n+1$, since
$2^{d_n}\mid 3m_n+1$. \status{proved; formalized, \S\ref{sec:lean}}
\end{proposition}

\begin{proof}
For the cap, $2^{d_n}m_{n+1}=3m_n+1$ with $m_{n+1}\ge1$. For the sufficient condition, multiply the
target $3^{n+1}<2^{S_{n+1}}$ by $m_n$: by Lemma~\ref{lem:aggregate},
$2^{S_{n+1}}m_n=2^{d}\,2^{S_n}m_n\ge 2^{d}3^nm_0>3\cdot3^nm_n=3^{n+1}m_n$ whenever
$2^dm_0>3m_n$.
\end{proof}

The cap therefore exceeds the sufficient threshold $3m_n/m_0$ by a factor of about $m_0$. The
underlying reason is the identity $\log_2 m_n=\log_2 m_0-R_n+E_n$ of Lemma~\ref{lem:ER}: sustained
negative drift makes the ordinary iterate \emph{larger}, and the ordinary height so created is
exactly the resource that permits the large valuations needed to repay the deficit.

\begin{remark}\label{rem:selffin}
What this rules out is a class of arguments, not a theorem: any argument of the form ``a
sufficiently deep drift excursion cannot recover, because the required valuation exceeds what the
current iterate can supply'' is blocked before it starts. Over all odd $m<200001$, the minimum over
all steps of $\log_2(3m_n+1)-(\al-R_n)$ --- available valuation minus what a one-step return to
$R=0$ requires --- is $+1.7370$, attained at the seed step of $m_0=3$, and is never negative.
(Revision~6 reported $+1.8301$; that is the minimum over steps $n\ge1$ only, and silently omitted
the seed step. The qualitative conclusion is unaffected.) \status{the inequality is proved; the margin is a
computational diagnostic}
\end{remark}

\section{Boundaries of the programme}\label{sec:boundaries}

This section collects the strongest no-go results. Their scope is deliberately narrow: each says
that a specific class of information cannot distinguish positive-integer realization. None says
that Collatz or EOC is unprovable.

\subsection{Finite cylinders do not encode infinite realization}\label{sec:cylinder}

\begin{proposition}\label{prop:cylinder}
For $k\ge1$, every odd residue class mod $2^k$ contains arbitrarily large positive odd integers.
Consequently every exact-realizer cylinder is nonempty and unbounded, and no property of integers
determined by finitely many $2$-adic digits can be equivalent to positivity.
\status{proved; formalized, \S\ref{sec:lean}}
\end{proposition}

\begin{proof}
$r+2^k(B+1)$ lies in the class and exceeds $B$; oddness is preserved since $2\mid 2^k$. For the
second claim, $1$ and $1-2^k$ share a class mod $2^k$ and differ in sign.
\end{proof}

The consequence is a quantifier distinction that should be displayed:
\[
\boxed{\ \forall N\ \exists m_N\ \text{(free)}\qquad\text{versus}\qquad \exists m\ \forall N\ \text{(the whole problem)}\ }
\]
A valuation prefix of length $N$ pins exactly one odd class mod $2^{S_N+1}$, which contains
unboundedly many positive integers; so finite-prefix realizability never obstructs. The content
lies entirely in whether \emph{one} integer works for all $N$, which is the bounded-prefix-realizer
equivalence of the companion work.

\begin{remark}[Scope]
This does \emph{not} say finite information is useless. Finite data can still prove quantitative
lower bounds --- Open Problem~C is precisely such a statement at each finite $N$. The claim is only
about \emph{deciding} the infinite realization property. \status{scope statement}
\end{remark}

\subsection{Coordinate sufficiency is not independent information}\label{sec:collapse}

Several coordinates introduced in the programme turn out to be invertible reparameterizations of
the exact realizer. The ledger below records each, with the reference where the collapse is proved.

\begin{center}
\small
\begin{tabular}{@{}p{0.30\textwidth}p{0.62\textwidth}@{}}
\toprule
Coordinate & Status \\
\midrule
Normalized carry $Z_N$ & Same moving-anchor information as $-C_N3^{-N}$; no independent obstruction. \\
Lift digits $t_N$ & Exactly the binary digits of a fixed realizer, read in blocks cut at the valuation sums. \\
Eventually-zero lift tail & Equivalent to eventual constancy of the prefix realizers, hence to positive-integer realization. \\
Terminal zero gap $g(D)$ & $\lceil \Edep(D)\rceil$; a restatement of the target, not a mechanism (\S\ref{sec:gap}). \\
Transport depths $F,M,H$ & Exact bookkeeping. On a genuine orbit the matching precision $M_j$ is the length of the zero-run of the \emph{seed}'s binary expansion beginning at bit $S_j$. \\
Cumulative excess regeneration & Bit intervals at distinct failures are disjoint, so the total is at most $\lfloor\log_2 m_0\rfloor+1$. \\
$B_N=C_N+3^N\chi$ & Equals $2^{S_N}m_N$ exactly for a realizing $\chi$; hence $v_2(B_N)=S_N$ exactly and $\log_2|B_N|-v_2(B_N)=\log_2 m_N$. \\
\bottomrule
\end{tabular}
\end{center}

The last row deserves emphasis. The Archimedean size and the $2$-adic valuation of $B_N$ are
\emph{complementary factors of one exact identity}, not two independent constraints; knowing the
word and either one determines the other. This is why attempts to intersect a $p$-adic condition
with an Archimedean one at this object produced nothing: there was only ever one condition. It also
explains why the product formula yields exactly the trivial bound $v_2(B)\le\log_2|B|$ here, and
why the odd-prime support of $B_N$ --- being that of the arbitrary odd integer $m_N$ --- is
unconstrained. \status{proved in companion audits; identities verified computationally without
exception}

\subsection{The known divergent-orbit word shadows do not characterize positive realization}\label{sec:cardinality}

\begin{theorem}[Word-shadow cardinality barrier]\label{thm:cardinality}
For $b\subseteq\mathbb{N}$ define the valuation word (digits indexed from $0$)
\[
w_b(i)=\begin{cases} 2 & \text{if } i\equiv2\ (\mathrm{mod}\ 3)\ \text{and}\ \lfloor i/3\rfloor\in b,\\
1 & \text{otherwise.}\end{cases}
\]
Then every $w_b$ is an admissible valuation word satisfying
\begin{enumerate}[label=(\alph*),leftmargin=2em]
\item $\sum_{i<n}w_b(i)\le n+\lfloor n/3\rfloor$, hence $2^{S_n}\le3^n$ for every $n$: zero
  confinement, in the integral form of Lemma~\ref{lem:integral};
\item $R_n\le(4/3-\al)n=-0.2516\ldots n\to-\infty$ linearly;
\item $\sum_n 2^{R_n}<\infty$;
\item $\#\{n:R_n\ge-G\}=O(G)$.
\end{enumerate}
The map $b\mapsto w_b$ is injective, so the family has cardinality $2^{\aleph_0}$. Since one
positive integer realizes at most one infinite valuation word, all but countably many members have
\emph{no} positive-integer realizer. \status{proved. Formalized: item (a) (as
\texttt{famWord\_sum\_le} with \texttt{famWord\_confined}), injectivity, and the abstract
countability argument. Items (b), (c), (d) are proved here but \emph{not} formalized, and the final
sentence is not instantiated at $w_b$ in Lean}
\end{theorem}

\begin{proof}
(a) At most one digit in three is a $2$, so $S_n\le n+\lfloor n/3\rfloor$; and
$2^{n+\lfloor n/3\rfloor}\le3^n$ follows by cubing both sides, since
$3(n+\lfloor n/3\rfloor)\le4n$ and $16\le27$. (b) $4/3<\al$. (c) The bound in (b) makes
$2^{R_n}$ dominated by a geometric series. (d) $R_n\ge-G$ forces $n\le G/0.2516$. Injectivity:
evaluate $w_b$ at $i=3k+2$. The last sentence is Cantor's theorem: a realizing assignment would
give an injection $\mathcal{P}(\mathbb{N})\to\mathbb{N}$.
\end{proof}

\begin{remark}[What this does and does not establish]\label{rem:cardscope}
The shared global word-level shadows established in Theorem~\ref{thm:cardinality} --- zero
confinement, linear negative drift, summability of $\sum2^{R_n}$, and the stated occupation bound
--- therefore do not characterize positive-integer realization. More generally, any proposed
criterion depending only on properties shared by this family cannot do so. This is a statement
about characterization, not about quantitative bounds: it does not preclude finite or word-level
information yielding realizer lower bounds. \status{scope statement}
\end{remark}

\begin{remark}[The barrier is sharply delimited]\label{rem:zcrecounter}
It is worth saying how far Theorem~\ref{thm:cardinality} reaches, because a broader reading has
circulated. Since the least prefix realizer $r(D)$ is a function of the valuation word alone,
\emph{bounded prefix realizers} is itself a global word-level property, and it characterizes
positive-integer realizability \emph{exactly}: a word has bounded prefix realizers if and only if
some positive odd integer realizes it (\S\ref{sec:cylinder}). It escapes the theorem because it is
not shared by the family $w_b$ --- it holds for the countably many realizable members and fails for
the rest.

So the barrier does not say that word-level information is insufficient in principle. It says that
the \emph{divergent-orbit shadows} --- the properties a hypothetical Type-II orbit's word is known
to have --- are too coarse, while the one word-level invariant sharp enough to decide the question
is a restatement of realizability rather than an independent handle. \status{Revision 7, scope}
\end{remark}

\begin{remark}[Explicit small witnesses]
The barrier is not an artefact of uncountability. The all-ones word (realizer $-1$) and the word
$1,2,1,2,\dots$ (realizer $-5$) are both zero-confined with $R_n\to-\infty$ linearly, and neither
has a positive-integer realizer. \status{computational verification of an exact statement}
\end{remark}

\subsection{Population sparsity is not pointwise placement}

Theorem~\ref{thm:GTC} and the entropy constant of \S\ref{sec:entropy} are population statements:
they count, respectively, orbit values in a window and words in a shell. Open Problem~C is a
statement about the \emph{minimum} over a family. A density-zero exceptional set may contain every
least-realizer champion, and no argument currently rules this out. The same gap defeats
density-one results in the literature (\S\ref{sec:prior}) as applied to launch integers, which by
construction form an exceptional set. \status{structural interpretation}

\subsection{Deep drift creates its own recovery capacity}

\S\ref{sec:selffin}. \status{proved}

\subsection{Single-window lifetime does not control occupation}

Remark~\ref{rem:LvsO}. \status{Revision 5, computational example}

\begin{remark}[Net effect]
These boundaries narrow the remaining search to actual positive-orbit arithmetic and to genuinely
pointwise or cross-scale mechanisms. They do not constitute an argument that the programme cannot
succeed.
\end{remark}

\section{The programme hierarchy}\label{sec:hierarchy}

\begin{center}
\small
\begin{tabular}{@{}llp{0.42\textwidth}@{}}
\toprule
Level & Statement & Status \\
\midrule
0 & Finite word realizability & Solved exactly (Prop.~\ref{prop:congruence}); every cylinder is
    nonempty and unbounded (Prop.~\ref{prop:cylinder}). \\
1 & Universal drift exit $(\mathrm{DE})$ & Equivalent to excluding divergent orbits
    (Thm.~\ref{thm:DEequiv}); compatible with nontrivial cycles; \emph{open}. \\
2 & Quantitative single-window: bounds on $L_c(m)$ & Transfers to frontier bounds
    (Prop.~\ref{prop:transfer}); no unconditional bound known (Rem.~\ref{sec:nounconditional}). \\
3 & Exponential single-window floor & Open Problem~\ref{op:C}; equivalent to
    $L_c(m)=O(\log m)$ (Prop.~\ref{prop:window}). \\
4 & Total occupation $\Occ_c(m)=O(\log m)$ & Full EOC (Conjecture~\ref{conj:exist}); implies
    Level~3, and implied by its corridor-uniform form \emph{given cycle exclusion}; strictness is
    \emph{open} (Rem.~\ref{rem:strictness}). \\
5 & Effective EOC $+$ finite cycle verification & Collatz (Thm.~\ref{thm:reduction},
    Cor.~\ref{cor:explicit}). \\
\bottomrule
\end{tabular}
\end{center}

Levels $5\Rightarrow4\Rightarrow3\Rightarrow2$ and $3\Rightarrow1$ hold; $2\Rightarrow1$ holds by
Remark~\ref{sec:nounconditional}. \emph{No converse is claimed}, and the reader should note what
that does and does not mean.

\begin{remark}[Strictness is not established by the pointwise gap]\label{rem:strictness}
Revision~6 asserted that ``no converse holds in general'' and that Level~4 is \emph{strictly}
stronger than Level~3. Those assertions are withdrawn as stated. The evidence offered for them was
the pointwise gap $\Occ_c>L_c$ (Rem.~\ref{rem:LvsO}), and that establishes only that the two
\emph{quantities} differ; it does not show that the two \emph{universal asymptotic bounds}
``$L_c(m)=O(\log m)$ for all $m$'' and ``$\Occ_c(m)=O(\log m)$ for all $m$'' are logically
inequivalent. Proving inequivalence would require exhibiting a consistent state of affairs in which
every single-window lifetime is logarithmic while some total occupation is not, and no such
argument is given here or known to us.

Moreover the evidence points the other way for one natural form of Level~3. If the lifetime bound
is assumed \emph{corridor-uniformly} --- for every $c'\ge0$ simultaneously, not at one fixed
corridor --- then Corollary~\ref{cor:Uall} bounds the injective initial segment, and
Proposition~\ref{prop:occstop} converts that into Level~4 \emph{given cycle exclusion}
(Rem.~\ref{rem:Uallscope}(i)). So corridor-uniform Level~3 implies Level~4 under that additional
assumption, and for that form strictness is at best unproved rather than established.

What survives: Level~4 implies Level~3 (immediately, since $L_c\le\Occ_c$); no implication from
Level~1 to cycle-freeness is established here, so the drift-exit route alone does not reach Collatz
(Cor.~\ref{cor:cycles}, Rem.~\ref{rem:strictlyweaker}) --- while \emph{neither} strictness
\emph{nor} non-entailment is claimed; and for a \emph{fixed} corridor the converse
Level~$3\Rightarrow$~Level~4 is open. \status{Revision 7, correction of scope}
\end{remark}

\section{The programme map}\label{sec:map}

\begin{center}
\scriptsize
\begin{tabular}{@{}p{0.085\textwidth}p{0.125\textwidth}p{0.155\textwidth}p{0.185\textwidth}p{0.145\textwidth}p{0.18\textwidth}@{}}
\toprule
Axis & Object & Exact bridge & Strongest proved & Open target & Closed routes \\
\midrule
Occupation & $\Occ_c(m)$ & Thm.~\ref{thm:reduction} to Collatz & Reduction to a finite cycle
search, given effective constants & Separated-return control & --- \\
Time & $L_c(m)$, $\tau_0(m)$ & Record chronology (Prop.~\ref{prop:inversion}); transfer
(Prop.~\ref{prop:transfer}) & $(\mathrm{DE})\iff$ no divergence (Thm.~\ref{thm:DEequiv}) &
Any unconditional bound on $L_c$ & Self-financing recovery blocks excursion arguments
(\S\ref{sec:selffin}) \\
Residue & $\rmin(N,c)$, $r(D)$ & Exact congruence (Prop.~\ref{prop:congruence}) + inversion &
Fixed-anchor sectors (Thms.~\ref{thm:periodic}--\ref{thm:evperiodic},
Prop.~\ref{prop:bulkdefect}) & Arbitrary irregular floor & Coordinate collapse
(\S\ref{sec:collapse}) \\
Divergence & $R_n$, $Q_n$, value set & Curry sparsity; last-maximum reduction
(Prop.~\ref{prop:lastmax}) & Reciprocal summability; $B<1/\beta_*$ (Thm.~\ref{thm:betastar}) &
Exclude divergence & The \emph{known} shadows are insufficient (Thm.~\ref{thm:cardinality}) \\
Symbolic & Valuation words; return labels & Realizer correspondence & Finite and infinite symbolic
freedom & Positive realization & Cardinality barrier (Thm.~\ref{thm:cardinality}) \\
\bottomrule
\end{tabular}
\end{center}

\subsection{Symbolic freedom, finite and infinite}\label{sec:symbolic}

Revision~5 established that no \emph{finite} forbidden-pattern mechanism can settle pointwise
balance for a natural mod-$32$ return observable: every finite binary history is realizable by a
genuine orbit, with prescribed events consecutive, and this was formally verified. The symbolic
conclusion is now stronger. A machine-checked construction produces an infinite zero-confined
valuation sequence carrying arbitrary prescribed binary labels at designated event positions: a
\emph{full shift} at the word level, inside the zero corridor.

The distinction that matters is therefore sharper than ``no finite forbidden patterns'': even
substantial \emph{infinite} symbolic freedom coexists with zero confinement. Positive-integer
realization of such words remains open and, by Theorem~\ref{thm:cardinality}, is not decided by the
global word-level shadows established there. The missing restriction is not local symbolic
dynamics.
\status{finite universality: Revision 5, formalized; infinite full shift: companion formalization}

\section{Closed routes}\label{sec:closed}

For each mechanism, what it achieved and where it stops. Status codes: \textsf{T} formal
result, \textsf{C} computational diagnostic, \textsf{I} interpretive closure.

\begin{center}
\scriptsize
\begin{tabular}{@{}rp{0.30\textwidth}p{0.48\textwidth}c@{}}
\toprule
& Mechanism & Where it stops & \\
\midrule
1 & Fixed-modulus drift packing & Bound $8/9$ proved; \emph{terminality} of the method rests on a
structural argument plus a finite scan (Rem.~\ref{rem:terminal}); surpassed by a different
mechanism & \textsf{T}/\textsf{C} \\
2 & Bulk entropy / $I_{\mathrm{Collatz}}$ & Population scale only; no pointwise
anti-concentration & \textsf{I} \\
3 & Shell injectivity & Counts, does not place & \textsf{I} \\
4 & Moving-anchor collision law & Symbolic prefix geometry only & \textsf{I} \\
5 & Chang finite/local exclusions & Finite \emph{and} infinite symbolic freedom
(\S\ref{sec:symbolic}) & \textsf{T} \\
6 & Lift digits & Binary blocks of the realizer (\S\ref{sec:collapse}) & \textsf{T} \\
7 & Terminal zero gap & $\lceil\Edep\rceil$: target, not mechanism (\S\ref{sec:gap}) & \textsf{T} \\
8 & Normalized carry & Same anchor information & \textsf{T} \\
9 & Transport survival budget & Exact bookkeeping; on genuine orbits matching precision collapses
to seed bits, cumulative excess $O(\log m_0)$ & \textsf{T} \\
10 & External $p$-adic arithmetic & No applicable uniform bound found; the dimension grows with
$N$ (\S\ref{sec:external}) & \textsf{I} \\
11 & The \emph{known divergent-orbit} word shadows & Cardinality barrier
(Thm.~\ref{thm:cardinality}); word-level criteria in general are \emph{not} excluded
(Rem.~\ref{rem:zcrecounter}) & \textsf{T} \\
12 & Corridor-uniform lifetime hypothesis & Equivalent in strength to $O(\log n)$ stopping
(Cor.~\ref{cor:Uall}) & \textsf{T} \\
13 & Deep-drift recovery impossibility & Self-financing (\S\ref{sec:selffin}) & \textsf{T} \\
\bottomrule
\end{tabular}
\end{center}

\subsection{External arithmetic}\label{sec:external}

We record only the durable conclusion. The natural cleared-denominator object attached to a
confined word is $B_N(D,\chi)=C_N+3^N\chi$, an $(N{+}1)$-term sum of $\{2,3\}$-units whose
$2$-adic valuation one would like to bound. Three families of deep results were examined against
it. $p$-adic lower bounds for linear forms in logarithms bound $\alpha_1^{b_1}\cdots\alpha_n^{b_n}-1$,
a \emph{product} minus one, never a sum of three or more terms; their constants are at best
geometric in the number of logarithms; and at $p=2$ with $2$ among the bases the hypotheses fail
and the quantity is identically zero. Quantitative results on $S$-unit equations bound the
\emph{number} of non-degenerate solutions, not their size, and are ineffective beyond two
variables; moreover explicit lower bounds for that count forbid any bound uniform in the number of
variables. The $p$-adic Subspace Theorem does apply formally and yields a bound of the right shape,
but it is ineffective, its subspace count and height threshold depend on the dimension, and here
the dimension is $N{+}1$. In every case the obstruction is the same: the number of additive terms
grows with $N$, while the available theorems are quantified ``for each fixed dimension'' with
constants that provably blow up. \status{applicability audit; external theorems cited but not
reproved}

\section{Open problems}\label{sec:open}

\begin{openproblem}[Effective EOC constants]\label{op:A}
Establish Conjecture~\ref{conj:eff} at any explicit $(c,K,C)$, or any explicit upper bound on
admissible constants in Conjecture~\ref{conj:exist}. This closes the Collatz implication outright
(Corollary~\ref{cor:explicit}).
\end{openproblem}

\begin{openproblem}[Cycle exclusion below threshold]\label{op:B}
Show no nontrivial accelerated cycle has minimum member below $m^*(K,C,c)$. Settled for
$m^*\le2^{71}$ by \cite{Barina2025}; open beyond.
\end{openproblem}

\begin{openproblem}[Exponential realizer floor / logarithmic lifetime]\label{op:C}
Find $\varepsilon>0$ with, equivalently (Proposition~\ref{prop:window}),
\[
\text{(residue) } \rmin(N,c)\ \ge\ 2^{\varepsilon N-O(1)}
\qquad\Longleftrightarrow\qquad
\text{(time) } L_c(m)\ \le\ \varepsilon^{-1}\log_2 m+O(1),
\]
for all relevant $c$-confined words. This is the residue--time core of EOC. Revision~5's separate
Open Problems C and D are the same question at two levels of precision and are merged here;
Revision~5's Open Problem~E (frontier defect) is its sharp-leading form, not an independent weaker
target (Remark~\ref{rem:OPE}).
\end{openproblem}

\begin{openproblem}[Intermediate time-axis bounds]\label{op:intermediate}
Prove \emph{any} unconditional pointwise upper bound on $L_c(m)$. By
Proposition~\ref{prop:transfer} even $L_c(m)\le Km$ yields a linear frontier, and
$L_c(m)\le K(\log_2 m)^2$ yields a stretched-exponential one. Nothing of this kind is known
(Remark~\ref{sec:nounconditional}), and any such bound would be new.
\end{openproblem}

\begin{openproblem}[Universal drift exit]\label{op:DE}
Show that for every positive odd $m$ there is $n\ge1$ with $3^n<2^{S_n(m)}$. By
Theorem~\ref{thm:DEequiv} this is equivalent to excluding divergent orbits; no implication from it
to cycle-freeness is established here (Rem.~\ref{rem:strictlyweaker}). This is the previously posed zero-corridor realizability frontier, reclassified as
the qualitative endpoint of the time axis (Remark~\ref{rem:DEprov}); it is not newly invented here.
\end{openproblem}

\begin{openproblem}[Separated returns]\label{op:F}
Control the \emph{weighted} total $\Occ_c=\sum_i\ell_i$ of the confined episodes. Revision~6 asked
for control of the episode count; the results below fix what such control can and cannot be.
\end{openproblem}

\begin{proposition}[Entry lemma]\label{prop:entry}
Let $a>0$ be a re-entry time, so $R_{a-1}>c\ge R_a$. Then $d_{a-1}=1$, and consequently
\[
0\ \le\ c-R_a\ <\ \al-1=\log_2(3/2)=0.58496\ldots
\]
\emph{Proof.} $R_a=R_{a-1}+d_{a-1}-\al$ with $1<\al<2$ and $d_{a-1}$ a positive integer; a digit
$\ge2$ would give $R_a>R_{a-1}>c$. \qed

So only the \emph{first} episode has threshold $c$: every later episode begins in a corridor of
width less than $0.585$, a bound independent of $c$, of $m_0$, and of how deep the drift fell
during the excursion. Drift may of course go far below $c$ \emph{inside} an episode.
\status{Revision 7; formalized, \S\ref{sec:lean}}
\end{proposition}

\begin{proposition}[Episode counts are unbounded]\label{prop:episodes}
For every $c\ge0$ and every $B$ there is a positive odd $m_0$ whose genuine accelerated orbit has
more than $B$ corridor re-entries before some finite horizon $N$. One may additionally take $m_0$
larger than any prescribed bound.

\emph{Proof sketch.} The word playing digit~$2$ inside the corridor and digit~$1$ outside has
drift confined to a fixed two-sided window, re-enters immediately after each exit, and exits again
within three steps; every finite prefix is realized by a positive odd integer via
Proposition~\ref{prop:congruence}, and the realizers of a prefix form a full residue class.
\qed
\status{Revision 7; formalized, \S\ref{sec:lean}}
\end{proposition}

\begin{remark}[What this rules out, and how far]\label{rem:episodescope}
A \emph{uniformly bounded} episode count is therefore false, and cannot be assumed as a closing
hypothesis. Two scope points must travel with this.

\emph{Finite prefix versus total occupation.} Propositions~\ref{prop:entry} and
\ref{prop:episodes} and everything in the next paragraph concern the realized prefix $[0,N)$.
Nothing is claimed about the \emph{total} occupation of the constructed seeds, which depends on the
orbit beyond $N$.

\emph{The coarse per-episode accounting is lossy, by a stated factor.} Under a single-window bound
each episode is charged $\ell_i\le1+K(\log_2m_{a_i}+c_i)$, and every such summand is at least
$\log_2m_0-c$. Along the family of Proposition~\ref{prop:episodes} the prefix carries at least $B$
re-entries while the seed may be taken with $m_0<2^{6B+O_c(1)}$, so the surrogate total
$Q\ge(B+1)(\log_2m_0-c)$ while the prefix occupation it bounds is at most $O_c(1)+2B$; hence
\[
Q/\Occ_c^{\mathrm{prefix}}\ \ge\ (\log_2m_0-c)/(2+O_c(1))\ =\ \Omega(\log m_0)
\]
along an unbounded sequence of realizing seeds. The per-episode single-window accounting therefore
overestimates by a factor $\Omega(\log m_0)$, and no choice of constants repairs it. This is an
obstruction to \emph{that accounting route} only; it is not an impossibility theorem about EOC, and
says nothing about arguments that do not decompose per episode.

What the problem needs, then, is not a bound on the episode count but a bound on $\sum_i\ell_i$
that does not factor through per-episode single-window estimates. \status{Revision 7}
\end{remark}

\begin{openproblem}[Density-one EOC]\label{op:G}
Can a density-one form of EOC be proved by the methods of Tao \cite{Tao}? This remains meaningful:
divergent-orbit sparsity says nothing about the occupation of orbits that do converge.
\end{openproblem}

\begin{openproblem}[Pointwise return balance]\label{op:H}
For an infinite aperiodic orbit, does the natural mod-$32$ return label equidistribute along the
orbit? By \S\ref{sec:symbolic} no finite --- and no word-level infinite --- symbolic restriction
settles this. Demoted here from a top-level EOC problem to an auxiliary one.
\end{openproblem}

\section{Formal verification status}\label{sec:lean}

The accompanying Lean~4 development (Mathlib) formalizes the elementary bridges of this paper. The
following are machine-checked, with no \texttt{sorry}, \texttt{admit}, added axiom, or
\texttt{opaque} escape; each depends only on the standard axioms \texttt{propext},
\texttt{Classical.choice}, \texttt{Quot.sound} (several on a strict subset).

Two procedural notes, added in Revision~7 after an audit of this table. First, every module named
below is now imported by the library root, so the default build target checks them; previously all
but one were reachable only by building them individually by name, and ``machine-checked'' was not
evidenced by the repository's default target. Second, the axiom claim is established by
\texttt{\#print axioms} run on each theorem; those commands are not themselves committed, so the
claim is reproducible but not self-evidencing from the source tree. Descriptions below name the
hypotheses each formalization actually carries --- several are conditional on
\texttt{ReciprocalSummable}, which is Theorem~\ref{thm:GTC}'s conclusion and is \emph{not}
formalized.

\begin{center}
\small
\begin{tabular}{@{}p{0.25\textwidth}p{0.16\textwidth}p{0.535\textwidth}@{}}
\toprule
Module & Paper reference & Content \\
\midrule
\texttt{CurryFoundation} & Prop.~\ref{prop:lastmax} & Curry identities; last-global-maximum
existence (unconditional); zero-corridor tail and deficit results, these under
\texttt{ReciprocalSummable}. \\
\texttt{ChangHistory},\\ \texttt{ChangFullShift} & \S\ref{sec:symbolic} & Finite return-history
universality; infinite zero-corridor full shift. \\
\texttt{ZCRERealizerGrowth} & \S\ref{sec:cylinder} & Bounded prefix realizers $\iff$
positive-integer realizability. \\
\texttt{RealizerLiftDigit},\\ \texttt{TerminalZeroGap} & \S\ref{sec:collapse}, \S\ref{sec:gap} &
Lift digits are binary blocks; the discrete gap. (The bracket $g=\lceil\Edep\rceil$ is in
\texttt{TransportRegeneration}, not here.) \\
\texttt{TransportRegen\-eration},\\ \texttt{SuperBudgetCan\-cellation} & \S\ref{sec:collapse} &
Transport bookkeeping; the gap bracket; disjoint-run bound (disjointness is a hypothesis, not a
conclusion). The forced/excess split is \emph{not} formalized. \\
\texttt{ConstraintHeight} & \S\ref{sec:height} & Height inequality and the separation criterion. \\
\texttt{GlobalSeparation} & Prop.~\ref{prop:cylinder}, Thm.~\ref{thm:cardinality} & Cylinder
positivity; finite-precision sign impossibility; the family $w_b$, its confinement and
injectivity; the countability barrier. \\
\texttt{OrbitLifetime} & Lem.~\ref{lem:aggregate}, Prop.~\ref{prop:corridorauto} & Aggregate
identity; growth envelope; automatic corridor. \\
\texttt{DriftExit} & Props.~\ref{prop:futuremin}, \ref{prop:selffin}, \ref{prop:transfer} &
Future-minimum theorem; cycle drift; self-financing inequality; frontier inversion. \\
\texttt{SeparatedReturns} & Props.~\ref{prop:entry}, \ref{prop:corridorvalue},
\ref{prop:episodes} & Restart relation; entry lemma; the corridor value criterion
(\texttt{corridor\_value\_lower}); occupation bounded by the arrival time, \emph{assuming} arrival
at $1$; the oscillating word and unbounded realizable episode counts. \\
\texttt{ZeroConfinedSeed} & Prop.~\ref{prop:lastmax} & The restart at the last drift maximum, as an
\emph{actual positive odd seed}: zero-confined at every horizon, with drift $\to-\infty$ ---
\emph{under the hypothesis} \texttt{ReciprocalSummable}, i.e. conditional on the external
Theorem~\ref{thm:GTC}. \\
\texttt{CurryInterface} & \S\ref{sec:divergence}, Thm.~\ref{thm:DEequiv} & Injective $\iff$
divergent; the external inputs as explicit hypotheses; divergence $\Rightarrow$ a zero-confined
positive seed; universal drift exit $\Rightarrow$ no divergent orbit. \\
\texttt{PowerSavingSummable} & Prop.~\ref{prop:recip} & The dyadic summation: a windowed
power-saving value count, with injectivity, yields reciprocal summability. \\
\bottomrule
\end{tabular}
\end{center}

\begin{remark}[The divergence chain, layer by layer]\label{rem:divledger}
Revision~6 listed reciprocal summability (Prop.~\ref{prop:recip}) as an unformalized external
input. It no longer is: \texttt{PowerSavingSummable} \emph{derives} it in Lean from the windowed
sparsity count, by the dyadic summation, with injectivity carrying the value-set sum to the
orbit-indexed one. The ledger is now
\[
\underbrace{\text{windowed sparsity}}_{\text{external theorem; audited, not formalized}}
\ +\
\underbrace{(\mathrm{DE})}_{\text{open hypothesis --- the target}}
\ \Longrightarrow\ \text{no divergent orbits},
\]
and both must be retained in any summary. The two are of different kinds: formalizing
Theorem~\ref{thm:GTC} would remove the \emph{external} dependency, and would neither discharge
$(\mathrm{DE})$ nor prove divergence exclusion. \status{Revision 7}
\end{remark}

Statements in this paper that are \emph{not} formalized include Theorem~\ref{thm:GTC} (external
input; audited from the source text), Corollary~\ref{cor:Uall}'s analytic absorption step
(Rem.~\ref{rem:Uallscope}(ii)), the proved-sector floors of \S\ref{sec:residue} (companion note),
and every conjecture.

\section{Computational calibration}\label{sec:comp}

We retain only computations that materially support interpretation, and label each as a
diagnostic.

\begin{itemize}[leftmargin=1.4em]
\item \emph{Frontier records.} A direct seed scan, justified by Proposition~\ref{prop:inversion},
  produces the least realizers of confined words at each depth. The maximum observed ratio of
  occupation to $\log_2$ seed remains under $10$ in the scanned range.
\item \emph{Identity verification.} The identities of Lemmas~\ref{lem:ER},~\ref{lem:aggregate} and
  \ref{lem:integral}, the future-minimum theorem, the descent criterion $m_n<m_0\iff R_n>E_n$, and
  the factorization $B_N=2^{S_N}m_N$ were each verified without exception over the ranges scanned.
\item \emph{Self-financing margin.} \S\ref{sec:selffin}, Remark~\ref{rem:selffin}.
\item \emph{Drift-exit diagnostics.} $\tau_0=\sigma$ was \emph{observed} for every odd seed tested
  away from cycle minima ($3\le m<60001$; $0$ exceptions), consistent with
  Remark~\ref{rem:howweak}. What is \emph{proved} there is only $\tau_0\le\sigma$ in general,
  together with the polynomial window to which any strict example is confined; the equality is
  not proved for all seeds.
\end{itemize}

\begin{remark}[Discipline]\label{rem:compdiscipline}
We make no empirical saturation claims. In particular we do not assert that any normalized
stopping-time or lifetime ratio is bounded: a comparison of a limited scan against published record
tables corrected an earlier such claim in the underlying audits. Where a ratio is quoted, the
scanned range is stated and no trend is inferred beyond it. Units also require care: the classical
total stopping time $\sigma_\infty$ counts \emph{even} steps --- equivalently the valuation depth
$S_L$ --- whereas $L$ counts \emph{odd} accelerated steps; the two must not be compared without
conversion. \status{methodological}
\end{remark}

\section{Relation to prior work}\label{sec:prior}

The occupation functional and its four tiers are, to our knowledge, not isolated elsewhere; the
Collatz reduction of Theorem~\ref{thm:reduction} is ours. Lemma~\ref{lem:ER} is due to
Eliahou~\cite{Eliahou} and Rozier~\cite{Rozier}. Theorem~\ref{thm:GTC} is Garcia--Tal
\cite{GarciaTal} with Curry's explicit form \cite{Curry}. The $2$-adic conjugacy underlying the
realizer correspondence is Bernstein \cite{Bernstein} and Bernstein--Lagarias
\cite{BernsteinLagarias}; cycle-structure bounds are Steiner \cite{Steiner}, Simons--de Weger
\cite{SimonsdeWeger} and Hercher \cite{Hercher}; verification is Ba\v{r}ina \cite{Barina2021,Barina2025};
stochastic models are Lagarias--Weiss \cite{LagariasWeiss} and Kontorovich--Lagarias
\cite{KontorovichLagarias}. Sturmian words are treated in Lothaire \cite{Lothaire}. The return
observable and its structural reduction are Chang's \cite{Chang1,Chang2}.

On the time axis specifically: the condition ``$\exists n:S_n>\al n$'' does not appear in the
surveyed literature under a standard name. It is not Terras's stopping time \cite{Terras}, which is
the condition $m_n<m_0$, i.e. $R_n>E_n$ (Proposition~\ref{prop:weaker}). The classical results on
stopping time --- Terras \cite{Terras}, Everett \cite{Everett}, and in a different framework Tao
\cite{Tao} --- are density statements and carry no pointwise content; none bounds the number of
steps an orbit may spend at or above its starting value in terms of that value, and such a bound
would prove Collatz by strong induction. \status{literature audit}

\section{Conclusion}\label{sec:conclusion}

EOC is best viewed as a programme with an exact residue--time duality and a global occupation
layer above it --- one that implies the single-window layer, and whose converse at a fixed corridor
is open rather than refuted (Rem.~\ref{rem:strictness}). The residue axis is solved for finite realizability and for several
fixed-anchor sectors, and remains open at the arbitrary-word frontier. The time axis separates into
a qualitative drift-exit problem --- equivalent, under the current divergence theory, to excluding
divergent orbits, with no implication to cycle-freeness established --- and quantitative lifetime
bounds that invert
directly to realizer floors on a graded scale. Garcia--Tal/Curry supplies strong unconditional
divergence-side sparsity and the last-maximum reduction, but not drift exit.

A sustained programme of structural audits has narrowed what can supply the missing information.
Finite cylinders cannot encode infinite realization; several apparently independent coordinates are
invertible reparameterizations of the exact realizer, and the Archimedean and $2$-adic readings of
the natural divisibility object are complementary factors of one identity; the \emph{known divergent-orbit}
valuation-word shadows cannot characterize positive realization, by an explicit uncountable family
--- though word-level criteria in general are not excluded, and one such criterion is exact
(Rem.~\ref{rem:zcrecounter}); launch
renormalization validly reaches actual positive-orbit dynamics, but its original corridor-uniform
closure is essentially an $O(\log n)$ stopping-time conjecture; and deep negative drift is
self-financing with respect to the recovery valuations it requires.

The surviving frontier is therefore actual positive-orbit arithmetic, at three nested levels:
qualitative exclusion of an infinite zero-corridor tail; quantitative control of single-window
lifetime, at any scale; and, beyond that, recurrence control sufficient for total occupation.
Neither EOC nor the Collatz conjecture is proved here.

\begin{thebibliography}{99}
\bibitem{Eliahou} S. Eliahou, \emph{The $3x+1$ problem: new lower bounds on nontrivial cycle
lengths}, Discrete Mathematics \textbf{118} (1993), 45--56.
\bibitem{Rozier} O. Rozier, \emph{The $3x+1$ problem: a lower bound hypothesis}, Functiones et
Approximatio Commentarii Mathematici \textbf{56} (2017), no.~1, 7--23.
\bibitem{Terras} R. Terras, \emph{A stopping time problem on the positive integers}, Acta
Arithmetica \textbf{30} (1976), 241--252.
\bibitem{Everett} C. J. Everett, \emph{Iteration of the number-theoretic function $f(2n)=n$,
$f(2n+1)=3n+2$}, Advances in Mathematics \textbf{25} (1977), 42--45.
\bibitem{GarciaTal} M. V. P. Garcia and F. A. Tal, \emph{A note on the generalized $3n+1$ problem},
Acta Arithmetica \textbf{90} (1999), no.~3, 245--250.
\bibitem{Curry} M. J. Curry, \emph{An Explicit Windowed Sparsity Bound for Divergent $3x+1$ Orbits,
with Logarithmic-Floor Exclusion beyond the Harmonic Barrier}, Zenodo, 24 August 2026.
\doi{10.5281/zenodo.22087163}
\bibitem{LagariasWeiss} J. C. Lagarias and A. Weiss, \emph{The $3x+1$ problem: two stochastic
models}, Annals of Applied Probability \textbf{2} (1992), 229--261.
\bibitem{KontorovichLagarias} A. V. Kontorovich and J. C. Lagarias, \emph{Stochastic models for the
$3x+1$ and $5x+1$ problems and related problems}, in: The Ultimate Challenge: The $3x+1$ Problem,
Amer. Math. Soc., 2010, 131--188.
\bibitem{Tao} T. Tao, \emph{Almost all orbits of the Collatz map attain almost bounded values},
Forum of Mathematics, Pi \textbf{10} (2022), e12.
\bibitem{Steiner} R. P. Steiner, \emph{A theorem on the Syracuse problem}, Proc. Seventh Manitoba
Conf. Numer. Math. and Computing, 1977, 553--559.
\bibitem{SimonsdeWeger} J. Simons and B. de Weger, \emph{Theoretical and computational bounds for
$m$-cycles of the $3n+1$ problem}, Acta Arithmetica \textbf{117} (2005), 51--70.
\bibitem{Hercher} C. Hercher, \emph{There are no Collatz $m$-cycles with $m\le91$}, Journal of
Integer Sequences \textbf{26} (2023), Article 23.3.5.
\bibitem{Barina2021} D. Ba\v{r}ina, \emph{Convergence verification of the Collatz problem}, The
Journal of Supercomputing \textbf{77} (2021), 2681--2688. \doi{10.1007/s11227-020-03368-x}
\bibitem{Barina2025} D. Ba\v{r}ina, \emph{Improved verification limit for the convergence of the
Collatz conjecture}, The Journal of Supercomputing \textbf{81} (2025), article~810.
\doi{10.1007/s11227-025-07337-0}
\bibitem{Lothaire} M. Lothaire, \emph{Algebraic Combinatorics on Words}, Cambridge University
Press, 2002. (Sturmian words: Chapter~2, by J. Berstel and P. S\'e\'ebold.)
\bibitem{DJperiodic} E. De Jes\'us, \emph{Periodic and Eventually Periodic Realizer Floors in the
Accelerated Collatz Carry Equation: Fixed $2$-Adic Anchors and Binary Periodicity}, Zenodo, 2026.
\doi{10.5281/zenodo.21889748}
\bibitem{DJentropy} E. De Jes\'us, \emph{The Entropy-Deficit Identity in the Accelerated Collatz
Word Model: First-Moment Closure, Occupation-Tail Rates, the Geometric Codimension, and the
Constant $0.0793186127\dots$}, Zenodo, 2026. \doi{10.5281/zenodo.21889866}
\bibitem{DJlaunch} E. De Jes\'us, \emph{Conditional Obstructions and Launch Renormalization in
Accelerated Collatz Dynamics}, companion note, 2026.
\bibitem{DJtransport} E. De Jes\'us, \emph{Transport Deficits and $2$-Adic Survival Budgets in
Accelerated Collatz Words}, companion note, 2026.
\bibitem{Bernstein} D. J. Bernstein, \emph{A non-iterative $2$-adic statement of the $3N+1$
conjecture}, Proc. Amer. Math. Soc. \textbf{121} (1994), 405--408.
\bibitem{BernsteinLagarias} D. J. Bernstein and J. C. Lagarias, \emph{The $3x+1$ conjugacy map},
Canadian Journal of Mathematics \textbf{48} (1996), 1154--1169.
\bibitem{Chang1} E. Y. Chang, \emph{A Structural Reduction of the Collatz Conjecture to One-Bit
Orbit Mixing}, preprint.
\bibitem{Chang2} E. Y. Chang, \emph{Exploring Collatz Dynamics with Human--LLM Collaboration},
preprint.
\end{thebibliography}

\appendix

\section{Band rigidity}\label{sec:band}

We retain the following structural observation from Revision~5, which constrains how a drift path
may sit inside a unit band.

\begin{theorem}[Band rigidity]\label{thm:band}
Let $\theta=\al-1$. If a word's drift path satisfies $\beta\le R_j<\beta+1$ for all $j$ in a window,
then the valuation letters in that window are determined up to the Beatty structure of $\theta$: the
admissible letters are confined to $\{1,2\}$ and their placement follows the Beatty gap sequence of
$\al$. \status{Revision 5, retained}
\end{theorem}

The relevance to this revision is limited but worth recording: band rigidity says that a
\emph{narrow} corridor forces near-Sturmian structure, which is exactly the regime where
Observation~\ref{obs:height} shows height amortization fails. Narrow corridors are therefore not
automatically easier.

\section{Claim-status conventions}\label{sec:status}

Every substantive statement carries a bracketed status:
\textsf{[Revision 5]} retained unchanged from the previous revision;
\textsf{[proved]} proved in this paper;
\textsf{[companion work]} proved in a cited companion note;
\textsf{[external, cited]} a cited theorem of others, used but not reproved;
\textsf{[formalized]} machine-checked, see \S\ref{sec:lean};
\textsf{[computational diagnostic]} supported by computation only, never promoted to a theorem;
\textsf{[structural interpretation]} an organizing reading of proved facts;
\textsf{[reclassification]} a previously known statement placed in a new role;
\textsf{[correction to Revision 5]} a superseded claim, corrected here.

\end{document}
